\documentclass[12pt]{article}
\usepackage[a4paper,margin=1in]{geometry}
\usepackage{iftex}
\ifXeTeX
\usepackage{fontspec}
\defaultfontfeatures{Ligatures=TeX}
\setmainfont{Times New Roman}
\else
\usepackage[T1]{fontenc}
\usepackage[utf8]{inputenc}
\usepackage{newtxtext}
\usepackage{newtxmath}
\fi
\usepackage{enumitem}
\usepackage{hyperref}
\usepackage{parskip}
\usepackage{microtype}
\usepackage[normalem]{ulem}
\usepackage{xcolor}

\newcommand{\practiceid}[1]{\texttt{#1}}
\newlength{\readinggutter}
\setlength{\readinggutter}{2.8em}
\newlength{\readingfirstindent}
\setlength{\readingfirstindent}{1.5em}
\newlength{\questionblockgap}
\setlength{\questionblockgap}{0.45\baselineskip}
\newlength{\exampleblockgap}
\setlength{\exampleblockgap}{0.65\baselineskip}
\newlength{\passageparagraphgap}
\setlength{\passageparagraphgap}{0.45\baselineskip}
\newenvironment{exampleblock}{\par\vspace{0.35\baselineskip}\begingroup\raggedright\setlength{\parskip}{0pt}\setlength{\parindent}{0pt}}{\par\endgroup\vspace{\exampleblockgap}}
\newcommand{\readinglineindent}[2]{%
\par\noindent%
\hangindent=\dimexpr\readinggutter+0.9em\relax%
\hangafter=1%
\makebox[\readinggutter][r]{\scriptsize\color{gray}#1}\hspace{0.9em}\hspace*{\readingfirstindent}#2%
\par}
\newcommand{\readingline}[2]{%
\par\noindent%
\hangindent=\dimexpr\readinggutter+0.9em\relax%
\hangafter=1%
\makebox[\readinggutter][r]{\scriptsize\color{gray}#1}\hspace{0.9em}#2%
\par}

\setlength{\parindent}{0pt}
\setlength{\parskip}{0.5em}
\setlist[enumerate]{topsep=0.25em,itemsep=0.18em,parsep=0pt,partopsep=0pt}
\setlength{\emergencystretch}{2em}

\begin{document}
\section*{TOEFL Practice 08 - Review}
Source of truth: DOCX only.
\section*{Listening Comprehension}
\subsection*{Part A}
DIRECTIONS: In Part A you will hear short conversations between two speakers. At the end of each conversation, a third speaker will ask a question about what the first two speakers said. Each conversation and each question will be spoken only one time. Therefore, you must listen carefully to understand what each speaker says. After you hear a conversation and the question, read the four choices and select the one that is the best answer to the question the speaker asked. Then, on your answer. sheet, find the number of the question and blacken the space that corresponds to the letter for the answer you have chosen. Blacken the space completely so that the letter inside the space does not show.\par
\begin{exampleblock}
Listen to the following example.\par
On the recording, you hear:\par
(Man) Does the car need to be filled?\par
(Woman) Mary stopped at the gas station on her way home.\par
(Narrator) What does the woman mean?\par
In your test book, you will read:\par
(A) Mary bought some food.\par
(B) Mary had car trouble.\par
(C) Mary went shopping.\par
(D) Mary bought some gas.\par
From the conversation you learn that Mary stopped at the gas station on her way home. The best answer to the question "Does the car need to be filled?" is (D), "Mary bought some gas." Therefore, the correct answer is (D).\par
Now let us begin Part A with question number 1.\par
\end{exampleblock}
\noindent\textbf{1.}
\begin{enumerate}[label=(\Alph*),leftmargin=*,itemsep=0.2em]
\item The soup is delicious.
\item The recipe is getting better.
\item The soup does not taste very good.
\item He decided to skip the soup.
\end{enumerate}
\par
\begingroup
\setlength{\parskip}{0pt}
\setlength{\parindent}{0pt}
\textbf{Review}\par
Answer: C\par
Confidence: medium\par
Jawaban diambil dari kunci belajar yang diberikan pengguna dan belum diverifikasi terhadap audio/transkrip resmi.\par
\endgroup
\par
\vspace{0.25\baselineskip}
\noindent\textbf{2.}
\begin{enumerate}[label=(\Alph*),leftmargin=*,itemsep=0.2em]
\item They have their troubles, and we have ours.
\item The man should write his own lab report.
\item They never tell us the correct hours.
\item The man needs to check the hours.
\end{enumerate}
\par
\begingroup
\setlength{\parskip}{0pt}
\setlength{\parindent}{0pt}
\textbf{Review}\par
Answer: D\par
Confidence: medium\par
Jawaban diambil dari kunci belajar yang diberikan pengguna dan belum diverifikasi terhadap audio/transkrip resmi.\par
\endgroup
\par
\vspace{0.25\baselineskip}
\noindent\textbf{3.}
\begin{enumerate}[label=(\Alph*),leftmargin=*,itemsep=0.2em]
\item Andy keeps track of advising and registration.
\item It is important to listen to Andy with skepticism.
\item It is not necessary to add salt to the dish.
\item Andy has already finished his registration.
\end{enumerate}
\par
\begingroup
\setlength{\parskip}{0pt}
\setlength{\parindent}{0pt}
\textbf{Review}\par
Answer: D\par
Confidence: medium\par
Jawaban diambil dari kunci belajar yang diberikan pengguna dan belum diverifikasi terhadap audio/transkrip resmi.\par
\endgroup
\par
\vspace{0.25\baselineskip}
\noindent\textbf{4.}
\begin{enumerate}[label=(\Alph*),leftmargin=*,itemsep=0.2em]
\item The lettuce needs to be rinsed before it is used.
\item Let us wash up before we eat the sandwiches.
\item We like eating lettuce because it is healthy.
\item The lettuce is in the salad she is making.
\end{enumerate}
\par
\begingroup
\setlength{\parskip}{0pt}
\setlength{\parindent}{0pt}
\textbf{Review}\par
Answer: A\par
Confidence: medium\par
Jawaban diambil dari kunci belajar yang diberikan pengguna dan belum diverifikasi terhadap audio/transkrip resmi.\par
\endgroup
\par
\vspace{0.25\baselineskip}
\noindent\textbf{5.}
\begin{enumerate}[label=(\Alph*),leftmargin=*,itemsep=0.2em]
\item At the School of Business.
\item At a bookstore.
\item In.a hardware shop.
\item In a grocery store.
\end{enumerate}
\par
\begingroup
\setlength{\parskip}{0pt}
\setlength{\parindent}{0pt}
\textbf{Review}\par
Answer: B\par
Confidence: medium\par
Jawaban diambil dari kunci belajar yang diberikan pengguna dan belum diverifikasi terhadap audio/transkrip resmi.\par
\endgroup
\par
\vspace{0.25\baselineskip}
\noindent\textbf{6.}
\begin{enumerate}[label=(\Alph*),leftmargin=*,itemsep=0.2em]
\item The sink has a bad odor.
\item It is too late to start looking at the puddle.
\item The man should not come at this hour.
\item The sink should be repaired.
\end{enumerate}
\par
\begingroup
\setlength{\parskip}{0pt}
\setlength{\parindent}{0pt}
\textbf{Review}\par
Answer: D\par
Confidence: medium\par
Jawaban diambil dari kunci belajar yang diberikan pengguna dan belum diverifikasi terhadap audio/transkrip resmi.\par
\endgroup
\par
\vspace{0.25\baselineskip}
\noindent\textbf{7.}
\begin{enumerate}[label=(\Alph*),leftmargin=*,itemsep=0.2em]
\item Mary is an author of a book.
\item Mary must have married a publisher.
\item The historian is getting married.
\item Historical novels are very popular.
\end{enumerate}
\par
\begingroup
\setlength{\parskip}{0pt}
\setlength{\parindent}{0pt}
\textbf{Review}\par
Answer: A\par
Confidence: medium\par
Jawaban diambil dari kunci belajar yang diberikan pengguna dan belum diverifikasi terhadap audio/transkrip resmi.\par
\endgroup
\par
\vspace{0.25\baselineskip}
\noindent\textbf{8.}
\begin{enumerate}[label=(\Alph*),leftmargin=*,itemsep=0.2em]
\item All customers get three pencils.
\item Those who take the test get a pencil.
\item All buyers receive a free gift.
\item Those who show their textbooks receive a free pencil.
\end{enumerate}
\par
\begingroup
\setlength{\parskip}{0pt}
\setlength{\parindent}{0pt}
\textbf{Review}\par
Answer: C\par
Confidence: medium\par
Jawaban diambil dari kunci belajar yang diberikan pengguna dan belum diverifikasi terhadap audio/transkrip resmi.\par
\endgroup
\par
\vspace{0.25\baselineskip}
\noindent\textbf{9.}
\begin{enumerate}[label=(\Alph*),leftmargin=*,itemsep=0.2em]
\item She has been looking for a cab for fifty minutes
\item She has been looking for a taxi for a pattern. quarter of an hour.
\item She has been lost for fifteen minutes while wearing a light coat.
\item She has sized up the situation correctly and called a cab
\end{enumerate}
\par
\begingroup
\setlength{\parskip}{0pt}
\setlength{\parindent}{0pt}
\textbf{Review}\par
Answer: B\par
Confidence: medium\par
Jawaban diambil dari kunci belajar yang diberikan pengguna dan belum diverifikasi terhadap audio/transkrip resmi.\par
\endgroup
\par
\vspace{0.25\baselineskip}
\noindent\textbf{10.}
\begin{enumerate}[label=(\Alph*),leftmargin=*,itemsep=0.2em]
\item People's recommendations are more effective than newspaper ads.
\item Most people do not pay attention to newspaper ads.
\item The new exhibit is promoted by the best advertising agency.
\item Words cannot describe the new exhibit effectively.
\end{enumerate}
\par
\begingroup
\setlength{\parskip}{0pt}
\setlength{\parindent}{0pt}
\textbf{Review}\par
Answer: A\par
Confidence: medium\par
Jawaban diambil dari kunci belajar yang diberikan pengguna dan belum diverifikasi terhadap audio/transkrip resmi.\par
\endgroup
\par
\vspace{0.25\baselineskip}
\noindent\textbf{11.}
\begin{enumerate}[label=(\Alph*),leftmargin=*,itemsep=0.2em]
\item The man needs to fill out the forms.
\item The man's car is out of gas and needs a fill-up.
\item The paint feels wet, and visitors should be careful.
\item The man is invited to go along.
\end{enumerate}
\par
\begingroup
\setlength{\parskip}{0pt}
\setlength{\parindent}{0pt}
\textbf{Review}\par
Answer: D\par
Confidence: medium\par
Jawaban diambil dari kunci belajar yang diberikan pengguna dan belum diverifikasi terhadap audio/transkrip resmi.\par
\endgroup
\par
\vspace{0.25\baselineskip}
\noindent\textbf{12.}
\begin{enumerate}[label=(\Alph*),leftmargin=*,itemsep=0.2em]
\item It is better to prepare before pouring it.
\item The batter contains four ingredients.
\item The mixture is poured on a hot surface.
\item The first pen is better than the next one.
\end{enumerate}
\par
\begingroup
\setlength{\parskip}{0pt}
\setlength{\parindent}{0pt}
\textbf{Review}\par
Answer: C\par
Confidence: medium\par
Jawaban diambil dari kunci belajar yang diberikan pengguna dan belum diverifikasi terhadap audio/transkrip resmi.\par
\endgroup
\par
\vspace{0.25\baselineskip}
\noindent\textbf{13.}
\begin{enumerate}[label=(\Alph*),leftmargin=*,itemsep=0.2em]
\item Julie and Michael plan to get married.
\item Julie and Michael decided to buy a new tie.
\item Julie's mother has not decided yet.
\item Julie's mother has no news to report.
\end{enumerate}
\par
\begingroup
\setlength{\parskip}{0pt}
\setlength{\parindent}{0pt}
\textbf{Review}\par
Answer: A\par
Confidence: medium\par
Jawaban diambil dari kunci belajar yang diberikan pengguna dan belum diverifikasi terhadap audio/transkrip resmi.\par
\endgroup
\par
\vspace{0.25\baselineskip}
\noindent\textbf{14.}
\begin{enumerate}[label=(\Alph*),leftmargin=*,itemsep=0.2em]
\item They cannot attend advising sessions.
\item An advisor is assigned to them.
\item They chose the wrong advising schedule.
\item An advisor shows them the correct schedule.
\end{enumerate}
\par
\begingroup
\setlength{\parskip}{0pt}
\setlength{\parindent}{0pt}
\textbf{Review}\par
Answer: B\par
Confidence: medium\par
Jawaban diambil dari kunci belajar yang diberikan pengguna dan belum diverifikasi terhadap audio/transkrip resmi.\par
\endgroup
\par
\vspace{0.25\baselineskip}
\noindent\textbf{15.}
\begin{enumerate}[label=(\Alph*),leftmargin=*,itemsep=0.2em]
\item It is a good idea to go jogging despitethe hot weather.
\item It is important to stand up straight andexercise regularly.
\item The current weather is difficult totolerate.
\item The weather has a long-standing pattern.
\end{enumerate}
\par
\begingroup
\setlength{\parskip}{0pt}
\setlength{\parindent}{0pt}
\textbf{Review}\par
Answer: C\par
Confidence: medium\par
Jawaban diambil dari kunci belajar yang diberikan pengguna dan belum diverifikasi terhadap audio/transkrip resmi.\par
\endgroup
\par
\vspace{0.25\baselineskip}
\noindent\textbf{16.}
\begin{enumerate}[label=(\Alph*),leftmargin=*,itemsep=0.2em]
\item They have little time for dancing these days.
\item They believe that gardening is serious business.
\item Their dancing classes have not started yet
\item Their gardening business is off to a good start.
\end{enumerate}
\par
\begingroup
\setlength{\parskip}{0pt}
\setlength{\parindent}{0pt}
\textbf{Review}\par
Answer: A\par
Confidence: medium\par
Jawaban diambil dari kunci belajar yang diberikan pengguna dan belum diverifikasi terhadap audio/transkrip resmi.\par
\endgroup
\par
\vspace{0.25\baselineskip}
\noindent\textbf{17.}
\begin{enumerate}[label=(\Alph*),leftmargin=*,itemsep=0.2em]
\item The color seems to be wrong.
\item She is not sure if she likes it.
\item She thinks that it is too short.
\item The man should not buy it.
\end{enumerate}
\par
\begingroup
\setlength{\parskip}{0pt}
\setlength{\parindent}{0pt}
\textbf{Review}\par
Answer: D\par
Confidence: medium\par
Jawaban diambil dari kunci belajar yang diberikan pengguna dan belum diverifikasi terhadap audio/transkrip resmi.\par
\endgroup
\par
\vspace{0.25\baselineskip}
\noindent\textbf{18.}
\begin{enumerate}[label=(\Alph*),leftmargin=*,itemsep=0.2em]
\item Renting a car
\item Going to the station
\item Planting a tree for children
\item Arriving late every day
\end{enumerate}
\par
\begingroup
\setlength{\parskip}{0pt}
\setlength{\parindent}{0pt}
\textbf{Review}\par
Answer: A\par
Confidence: medium\par
Jawaban diambil dari kunci belajar yang diberikan pengguna dan belum diverifikasi terhadap audio/transkrip resmi.\par
\endgroup
\par
\vspace{0.25\baselineskip}
\noindent\textbf{19.}
\begin{enumerate}[label=(\Alph*),leftmargin=*,itemsep=0.2em]
\item The search committee has filled out the applications.
\item The search for an office assistant has ended.
\item The hiring process has not been completed.
\item The applicants will still be reviewing the office
\end{enumerate}
\par
\begingroup
\setlength{\parskip}{0pt}
\setlength{\parindent}{0pt}
\textbf{Review}\par
Answer: C\par
Confidence: medium\par
Jawaban diambil dari kunci belajar yang diberikan pengguna dan belum diverifikasi terhadap audio/transkrip resmi.\par
\endgroup
\par
\vspace{0.25\baselineskip}
\noindent\textbf{20.}
\begin{enumerate}[label=(\Alph*),leftmargin=*,itemsep=0.2em]
\item Come back on Tuesday.
\item Take the test.
\item Take the course.
\item Make a choice of a major
\end{enumerate}
\par
\begingroup
\setlength{\parskip}{0pt}
\setlength{\parindent}{0pt}
\textbf{Review}\par
Answer: B\par
Confidence: medium\par
Jawaban diambil dari kunci belajar yang diberikan pengguna dan belum diverifikasi terhadap audio/transkrip resmi.\par
\endgroup
\par
\vspace{0.25\baselineskip}
\noindent\textbf{21.}
\begin{enumerate}[label=(\Alph*),leftmargin=*,itemsep=0.2em]
\item The dinner last week was a disaster.
\item She'll pay for the dinner.
\item She'll cook dinner this time.
\item She lives on Broadway.
\end{enumerate}
\par
\begingroup
\setlength{\parskip}{0pt}
\setlength{\parindent}{0pt}
\textbf{Review}\par
Answer: B\par
Confidence: medium\par
Jawaban diambil dari kunci belajar yang diberikan pengguna dan belum diverifikasi terhadap audio/transkrip resmi.\par
\endgroup
\par
\vspace{0.25\baselineskip}
\noindent\textbf{22.}
\begin{enumerate}[label=(\Alph*),leftmargin=*,itemsep=0.2em]
\item He finished the report already.
\item He wants to have another meeting.
\item He is tired' of meetings.
\item He doesn't have time on Sunday.
\end{enumerate}
\par
\begingroup
\setlength{\parskip}{0pt}
\setlength{\parindent}{0pt}
\textbf{Review}\par
Answer: C\par
Confidence: medium\par
Jawaban diambil dari kunci belajar yang diberikan pengguna dan belum diverifikasi terhadap audio/transkrip resmi.\par
\endgroup
\par
\vspace{0.25\baselineskip}
\noindent\textbf{23.}
\begin{enumerate}[label=(\Alph*),leftmargin=*,itemsep=0.2em]
\item Making an appointment
\item Going out on a date
\item Discussing Peter
\item Scheduling a departure
\end{enumerate}
\par
\begingroup
\setlength{\parskip}{0pt}
\setlength{\parindent}{0pt}
\textbf{Review}\par
Answer: A\par
Confidence: medium\par
Jawaban diambil dari kunci belajar yang diberikan pengguna dan belum diverifikasi terhadap audio/transkrip resmi.\par
\endgroup
\par
\vspace{0.25\baselineskip}
\noindent\textbf{24.}
\begin{enumerate}[label=(\Alph*),leftmargin=*,itemsep=0.2em]
\item He will rest in the afternoon.
\item He has to keep working.
\item He needs to clean his desk.
\item He has missed the deadline.
\end{enumerate}
\par
\begingroup
\setlength{\parskip}{0pt}
\setlength{\parindent}{0pt}
\textbf{Review}\par
Answer: B\par
Confidence: medium\par
Jawaban diambil dari kunci belajar yang diberikan pengguna dan belum diverifikasi terhadap audio/transkrip resmi.\par
\endgroup
\par
\vspace{0.25\baselineskip}
\noindent\textbf{25.}
\begin{enumerate}[label=(\Alph*),leftmargin=*,itemsep=0.2em]
\item Fred is mysterious.
\item Fred's class is not difficult.
\item She doesn't know.
\item She is taking a different class.
\end{enumerate}
\par
\begingroup
\setlength{\parskip}{0pt}
\setlength{\parindent}{0pt}
\textbf{Review}\par
Answer: C\par
Confidence: medium\par
Jawaban diambil dari kunci belajar yang diberikan pengguna dan belum diverifikasi terhadap audio/transkrip resmi.\par
\endgroup
\par
\vspace{0.25\baselineskip}
\noindent\textbf{26.}
\begin{enumerate}[label=(\Alph*),leftmargin=*,itemsep=0.2em]
\item Roberta had misplaced her key before.
\item Roberta's husband is late again.
\item Roberta's husband does not treat her very well.
\item Roberta had returned from work.
\end{enumerate}
\par
\begingroup
\setlength{\parskip}{0pt}
\setlength{\parindent}{0pt}
\textbf{Review}\par
Answer: A\par
Confidence: medium\par
Jawaban diambil dari kunci belajar yang diberikan pengguna dan belum diverifikasi terhadap audio/transkrip resmi.\par
\endgroup
\par
\vspace{0.25\baselineskip}
\noindent\textbf{27.}
\begin{enumerate}[label=(\Alph*),leftmargin=*,itemsep=0.2em]
\item At a gift shop
\item In a bank
\item At a florist's
\item In a delivery office
\end{enumerate}
\par
\begingroup
\setlength{\parskip}{0pt}
\setlength{\parindent}{0pt}
\textbf{Review}\par
Answer: C\par
Confidence: medium\par
Jawaban diambil dari kunci belajar yang diberikan pengguna dan belum diverifikasi terhadap audio/transkrip resmi.\par
\endgroup
\par
\vspace{0.25\baselineskip}
\noindent\textbf{28.}
\begin{enumerate}[label=(\Alph*),leftmargin=*,itemsep=0.2em]
\item It takes an hour to get to the NorthCampus from here.
\item The shuttle comes twice an hour.
\item It'll be 15 minutes for the shuttle to get here.
\item The shuttle is 15 minutes behind the schedule.
\end{enumerate}
\par
\begingroup
\setlength{\parskip}{0pt}
\setlength{\parindent}{0pt}
\textbf{Review}\par
Answer: B\par
Confidence: medium\par
Jawaban diambil dari kunci belajar yang diberikan pengguna dan belum diverifikasi terhadap audio/transkrip resmi.\par
\endgroup
\par
\vspace{0.25\baselineskip}
\noindent\textbf{29.}
\begin{enumerate}[label=(\Alph*),leftmargin=*,itemsep=0.2em]
\item The man can't park there.
\item It's okay to leave the car only for a few minutes.
\item She doesn't live on this block.
\item The man should park in the driveway.
\end{enumerate}
\par
\begingroup
\setlength{\parskip}{0pt}
\setlength{\parindent}{0pt}
\textbf{Review}\par
Answer: A\par
Confidence: medium\par
Jawaban diambil dari kunci belajar yang diberikan pengguna dan belum diverifikasi terhadap audio/transkrip resmi.\par
\endgroup
\par
\vspace{0.25\baselineskip}
\noindent\textbf{30.}
\begin{enumerate}[label=(\Alph*),leftmargin=*,itemsep=0.2em]
\item The woman is more careful nowthan she used to be.
\item The woman is not very careful.
\item He will help the woman with the gate.
\item He wishes only the best for the woman.
\end{enumerate}
\par
\begingroup
\setlength{\parskip}{0pt}
\setlength{\parindent}{0pt}
\textbf{Review}\par
Answer: B\par
Confidence: medium\par
Jawaban diambil dari kunci belajar yang diberikan pengguna dan belum diverifikasi terhadap audio/transkrip resmi.\par
\endgroup
\par
\vspace{0.25\baselineskip}
\subsection*{Part B}
DIRECTIONS: In this part of the test, you will hear longer conversations. After each conversation, you will hear several questions. The conversations and questions will not be repeated. After you hear a question, read the four possible answers in. your test book and select. The best answer. Then, on your answer sheet, find the number of the question and fill in the space that corresponds to the letter of the answer you have chosen. Remember, you are not allowed to take notes or write in your test book.\par
\begin{exampleblock}
Listen to the following example:\par
You will hear:\par
You will read:\par
(A) He has changed jobs.\par
(B) He has two children.\par
(C) He has two jobs.\par
(D) He is looking for a job.\par
From the conversation you learn that Tom has taken an additional job. The best answer to the question "Why is Tom tired?" is (C), "He has two jobs." Therefore the correct answer is (C).\par
\end{exampleblock}
\noindent\textbf{31.}
\begin{enumerate}[label=(\Alph*),leftmargin=*,itemsep=0.2em]
\item Select trees that he can plant in his yard.
\item Suggest trees that do not require much care.
\item Figure out the best variety of trees.
\item Point out trees that are on sale.
\end{enumerate}
\par
\begingroup
\setlength{\parskip}{0pt}
\setlength{\parindent}{0pt}
\textbf{Review}\par
Answer: B\par
Confidence: medium\par
Jawaban diambil dari kunci belajar yang diberikan pengguna dan belum diverifikasi terhadap audio/transkrip resmi.\par
\endgroup
\par
\vspace{0.25\baselineskip}
\noindent\textbf{32.}
\begin{enumerate}[label=(\Alph*),leftmargin=*,itemsep=0.2em]
\item Fruit growth, shade, and appearance
\item Soil erosion, height, and slope orientation
\item Southern exposure, roots, and water
\item Soil density, light, and heat tolerance
\end{enumerate}
\par
\begingroup
\setlength{\parskip}{0pt}
\setlength{\parindent}{0pt}
\textbf{Review}\par
Answer: D\par
Confidence: medium\par
Jawaban diambil dari kunci belajar yang diberikan pengguna dan belum diverifikasi terhadap audio/transkrip resmi.\par
\endgroup
\par
\vspace{0.25\baselineskip}
\noindent\textbf{33.}
\begin{enumerate}[label=(\Alph*),leftmargin=*,itemsep=0.2em]
\item They do not do well in direct sunlight.
\item Their root systems need good drainage.
\item They improve the environment and shade.
\item Their height can become a problem.
\end{enumerate}
\par
\begingroup
\setlength{\parskip}{0pt}
\setlength{\parindent}{0pt}
\textbf{Review}\par
Answer: B\par
Confidence: medium\par
Jawaban diambil dari kunci belajar yang diberikan pengguna dan belum diverifikasi terhadap audio/transkrip resmi.\par
\endgroup
\par
\vspace{0.25\baselineskip}
\noindent\textbf{34.}
\begin{enumerate}[label=(\Alph*),leftmargin=*,itemsep=0.2em]
\item Buy a few short trees.
\item Prune more trees.
\item Learn more about trees.
\item Think a little longer.
\end{enumerate}
\par
\begingroup
\setlength{\parskip}{0pt}
\setlength{\parindent}{0pt}
\textbf{Review}\par
Answer: C\par
Confidence: medium\par
Jawaban diambil dari kunci belajar yang diberikan pengguna dan belum diverifikasi terhadap audio/transkrip resmi.\par
\endgroup
\par
\vspace{0.25\baselineskip}
\noindent\textbf{35.}
\begin{enumerate}[label=(\Alph*),leftmargin=*,itemsep=0.2em]
\item Talk to the violin maker.
\item Play the instrument.
\item Wait for the wood to dry.
\item Adjust the bow strings.
\end{enumerate}
\par
\begingroup
\setlength{\parskip}{0pt}
\setlength{\parindent}{0pt}
\textbf{Review}\par
Answer: B\par
Confidence: medium\par
Jawaban diambil dari kunci belajar yang diberikan pengguna dan belum diverifikasi terhadap audio/transkrip resmi.\par
\endgroup
\par
\vspace{0.25\baselineskip}
\noindent\textbf{36.}
\begin{enumerate}[label=(\Alph*),leftmargin=*,itemsep=0.2em]
\item .It is a lengthy process.
\item It is a delightful hobby.
\item It is a musical endeavor.
\item It requires musical expertise.
\end{enumerate}
\par
\begingroup
\setlength{\parskip}{0pt}
\setlength{\parindent}{0pt}
\textbf{Review}\par
Answer: A\par
Confidence: medium\par
Jawaban diambil dari kunci belajar yang diberikan pengguna dan belum diverifikasi terhadap audio/transkrip resmi.\par
\endgroup
\par
\vspace{0.25\baselineskip}
\noindent\textbf{37.}
\begin{enumerate}[label=(\Alph*),leftmargin=*,itemsep=0.2em]
\item Music
\item Violin
\item Writing
\item Reading
\end{enumerate}
\par
\begingroup
\setlength{\parskip}{0pt}
\setlength{\parindent}{0pt}
\textbf{Review}\par
Answer: C\par
Confidence: medium\par
Jawaban diambil dari kunci belajar yang diberikan pengguna dan belum diverifikasi terhadap audio/transkrip resmi.\par
\endgroup
\par
\vspace{0.25\baselineskip}
\noindent\textbf{38.}
\begin{enumerate}[label=(\Alph*),leftmargin=*,itemsep=0.2em]
\item Glasses are less inconvenient.
\item Contacts can be more easily lost.
\item Contacts are more difficult to keepclean.
\item Glasses are less preferable.
\end{enumerate}
\par
\begingroup
\setlength{\parskip}{0pt}
\setlength{\parindent}{0pt}
\textbf{Review}\par
Answer: D\par
Confidence: medium\par
Jawaban diambil dari kunci belajar yang diberikan pengguna dan belum diverifikasi terhadap audio/transkrip resmi.\par
\endgroup
\par
\vspace{0.25\baselineskip}
\noindent\textbf{39.}
\begin{enumerate}[label=(\Alph*),leftmargin=*,itemsep=0.2em]
\item One
\item Two
\item Three
\item Four
\end{enumerate}
\par
\begingroup
\setlength{\parskip}{0pt}
\setlength{\parindent}{0pt}
\textbf{Review}\par
Answer: C\par
Confidence: medium\par
Jawaban diambil dari kunci belajar yang diberikan pengguna dan belum diverifikasi terhadap audio/transkrip resmi.\par
\endgroup
\par
\vspace{0.25\baselineskip}
\noindent\textbf{40.}
\begin{enumerate}[label=(\Alph*),leftmargin=*,itemsep=0.2em]
\item Researching lens plastics
\item Scheduling an eye exam
\item Locating an optometrist
\item Saving enough money
\end{enumerate}
\par
\begingroup
\setlength{\parskip}{0pt}
\setlength{\parindent}{0pt}
\textbf{Review}\par
Answer: B\par
Confidence: medium\par
Jawaban diambil dari kunci belajar yang diberikan pengguna dan belum diverifikasi terhadap audio/transkrip resmi.\par
\endgroup
\par
\vspace{0.25\baselineskip}
\subsection*{Part C}
DIRECTIONS: In Part C you will hear short lectures and extended conversations. At the end of each, you will be asked several questions. Each lecture or conversation and each question will be spoken only one time. For this reason, you must listen carefully to understand what each speaker says. After you hear a question, read the four choices and select the one that best answers the question the speaker asked. Then, on your answer sheet, find the number of the question and blacken the space that contains the letter for the answer you have chosen. Answer all questions according to what is stated or implied in the lecture or conversation.\par
\begin{exampleblock}
Listen to this sample talk.\par
You will hear:\par
Now listen, to the following example.\par
You will hear:\par
You will read:\par
(A) By cars and carriages\par
(B) By bicycles, trains, and carriages\par
(C) On foot and by boat\par
(D) On board ships and trains\par
The best answer to the question "According to the speaker, how did people travel before the invention of the automobile?" is (B), "By bicycles, trains, and carriages." Therefore, the correct answer is (B). You are not allowed to make notes during the test.\par
\end{exampleblock}
\noindent\textbf{41.}
\begin{enumerate}[label=(\Alph*),leftmargin=*,itemsep=0.2em]
\item To assist international travelers with data and changes
\item To elaborate on the system calendars in various locations
\item To explain the conventions of marking and changing dates
\item To demonstrate the connections between meridians and air travel
\end{enumerate}
\par
\begingroup
\setlength{\parskip}{0pt}
\setlength{\parindent}{0pt}
\textbf{Review}\par
Answer: C\par
Confidence: medium\par
Jawaban diambil dari kunci belajar yang diberikan pengguna dan belum diverifikasi terhadap audio/transkrip resmi.\par
\endgroup
\par
\vspace{0.25\baselineskip}
\noindent\textbf{42.}
\begin{enumerate}[label=(\Alph*),leftmargin=*,itemsep=0.2em]
\item Travelers have to choose the best days and times for their cross-continental journeys.
\item The distance from Los Angeles to Asia is similar to that from New York to London.
\item Date changes occur some time when air travelers cross the Pacific.
\item A loss of several hours of sleep is possible during long cross-continental flights.
\end{enumerate}
\par
\begingroup
\setlength{\parskip}{0pt}
\setlength{\parindent}{0pt}
\textbf{Review}\par
Answer: C\par
Confidence: medium\par
Jawaban diambil dari kunci belajar yang diberikan pengguna dan belum diverifikasi terhadap audio/transkrip resmi.\par
\endgroup
\par
\vspace{0.25\baselineskip}
\noindent\textbf{43.}
\begin{enumerate}[label=(\Alph*),leftmargin=*,itemsep=0.2em]
\item They should count the number of hours or add hours to the local time.
\item They need to add or subtract a day when crossing the date line.
\item They should buy a Western or Eastern calendar to match the local day count.
\item They need to count the beginnings of days from midnight to midnight.
\end{enumerate}
\par
\begingroup
\setlength{\parskip}{0pt}
\setlength{\parindent}{0pt}
\textbf{Review}\par
Answer: B\par
Confidence: medium\par
Jawaban diambil dari kunci belajar yang diberikan pengguna dan belum diverifikasi terhadap audio/transkrip resmi.\par
\endgroup
\par
\vspace{0.25\baselineskip}
\noindent\textbf{44.}
\begin{enumerate}[label=(\Alph*),leftmargin=*,itemsep=0.2em]
\item North-to-south standard time zones
\item Loose elliptical connections among longitudes
\item The line near the 24th hour on the local calendar
\item The area along the 180th meridian of longitude
\end{enumerate}
\par
\begingroup
\setlength{\parskip}{0pt}
\setlength{\parindent}{0pt}
\textbf{Review}\par
Answer: D\par
Confidence: medium\par
Jawaban diambil dari kunci belajar yang diberikan pengguna dan belum diverifikasi terhadap audio/transkrip resmi.\par
\endgroup
\par
\vspace{0.25\baselineskip}
\noindent\textbf{45.}
\begin{enumerate}[label=(\Alph*),leftmargin=*,itemsep=0.2em]
\item To describe the people who live on the island
\item To point out earth divisions similar to those in Alaska
\item To explain the rigidity of meridian lines
\item To illustrate the swerves in the date line
\end{enumerate}
\par
\begingroup
\setlength{\parskip}{0pt}
\setlength{\parindent}{0pt}
\textbf{Review}\par
Answer: D\par
Confidence: medium\par
Jawaban diambil dari kunci belajar yang diberikan pengguna dan belum diverifikasi terhadap audio/transkrip resmi.\par
\endgroup
\par
\vspace{0.25\baselineskip}
\noindent\textbf{46.}
\begin{enumerate}[label=(\Alph*),leftmargin=*,itemsep=0.2em]
\item The prehistoric northern settlements carbohydrates and the Northwestern Passage
\item The historical and economic significance of northern explorations
\item The Greek, English, and Spanish explorers of the northern seas
\item The connections between the English-Spanish war and northern explorers
\end{enumerate}
\par
\begingroup
\setlength{\parskip}{0pt}
\setlength{\parindent}{0pt}
\textbf{Review}\par
Answer: B\par
Confidence: medium\par
Jawaban diambil dari kunci belajar yang diberikan pengguna dan belum diverifikasi terhadap audio/transkrip resmi.\par
\endgroup
\par
\vspace{0.25\baselineskip}
\noindent\textbf{47.}
\begin{enumerate}[label=(\Alph*),leftmargin=*,itemsep=0.2em]
\item Together with the English, they defeated the Spanish fleet.
\item Unlike the Dutch, they looked for the Northern Passage.
\item Together with the Spanish, they controlled the southern seas.
\item Unlike other merchants, they did not trade in India.
\end{enumerate}
\par
\begingroup
\setlength{\parskip}{0pt}
\setlength{\parindent}{0pt}
\textbf{Review}\par
Answer: C\par
Confidence: medium\par
Jawaban diambil dari kunci belajar yang diberikan pengguna dan belum diverifikasi terhadap audio/transkrip resmi.\par
\endgroup
\par
\vspace{0.25\baselineskip}
\noindent\textbf{48.}
\begin{enumerate}[label=(\Alph*),leftmargin=*,itemsep=0.2em]
\item Those that are low in calories
\item Those that are high in protein
\item Those that are eaten on a usual basis
\item Those that are eaten at regular intervals
\end{enumerate}
\par
\begingroup
\setlength{\parskip}{0pt}
\setlength{\parindent}{0pt}
\textbf{Review}\par
Answer: C\par
Confidence: medium\par
Jawaban diambil dari kunci belajar yang diberikan pengguna dan belum diverifikasi terhadap audio/transkrip resmi.\par
\endgroup
\par
\vspace{0.25\baselineskip}
\noindent\textbf{49.}
\begin{enumerate}[label=(\Alph*),leftmargin=*,itemsep=0.2em]
\item Vitamins and minerals
\item Corn and potatoes
\item Bread and butter
\item Fish and beans
\end{enumerate}
\par
\begingroup
\setlength{\parskip}{0pt}
\setlength{\parindent}{0pt}
\textbf{Review}\par
Answer: D\par
Confidence: medium\par
Jawaban diambil dari kunci belajar yang diberikan pengguna dan belum diverifikasi terhadap audio/transkrip resmi.\par
\endgroup
\par
\vspace{0.25\baselineskip}
\noindent\textbf{50.}
\begin{enumerate}[label=(\Alph*),leftmargin=*,itemsep=0.2em]
\item Various kinds of fats
\item Different types of sugars
\item Half of all calories consumed by Americans
\item Half of all nutrients found in carbohydrates
\end{enumerate}
\par
\begingroup
\setlength{\parskip}{0pt}
\setlength{\parindent}{0pt}
\textbf{Review}\par
Answer: B\par
Confidence: medium\par
Jawaban diambil dari kunci belajar yang diberikan pengguna dan belum diverifikasi terhadap audio/transkrip resmi.\par
\endgroup
\par
\vspace{0.25\baselineskip}
\section*{Structure and Written Expression}
\subsection*{Sentence Completion}
This section is designed to test your ability to recognize language structures that are appropriate in standard written English. The questions in this section belong to two types, each of which has special directions.\par
DIRECTIONS: Questions 1-15 are partial sentences. Below each sentence you will see four words or phrases, marked (A), (B), (C), and (D). Select the one word or phrase that best completes the sentence. Then, on your answer sheet, find the number of the question and fill in the space that contains the letter for the answer you have chosen. Fill in the space completely.\par
\begin{exampleblock}
\textbf{EXAMPLE I}\par
\noindent Drying flowers is the best way......them.
\begin{enumerate}[label=(\Alph*),leftmargin=*,itemsep=0.2em]
\item to preserve
\item by preserving
\item preserve
\item preserved
\end{enumerate}
\vspace{0.25\baselineskip}
The sentence should state, "Drying flowers is the best way to preserve them." Therefore, the correct answer is (A).\par
\end{exampleblock}
\begin{exampleblock}
\textbf{EXAMPLE II}\par
\noindent Many American universities......as small, private colleges.
\begin{enumerate}[label=(\Alph*),leftmargin=*,itemsep=0.2em]
\item begun
\item beginning
\item began
\item for the beginning
\end{enumerate}
\vspace{0.25\baselineskip}
The sentence should state, "Many American universities began as small, private colleges." Therefore, the correct answer is (C). After you read the directions, begin work on the questions.\par
\end{exampleblock}
\noindent\textbf{1.} The Internet.... access to current news, political articles, business statistics, and software for practically any purpose.
\begin{enumerate}[label=(\Alph*),leftmargin=*,itemsep=0.2em]
\item can provide ready
\item can be provided readily
\item ready and can be providing
\item is ready and can provide
\end{enumerate}
\par
\begingroup
\setlength{\parskip}{0pt}
\setlength{\parindent}{0pt}
\textbf{Review}\par
Answer: A\par
Confidence: high\par
Grammar rule: Modal + base verb with object noun phrase.\par
Setelah modal `can`, diperlukan base verb. Pilihan (A) membentuk pola lengkap `can provide ready access`, dengan `access` sebagai objek noun phrase yang benar.\par
\endgroup
\par
\vspace{0.25\baselineskip}
\noindent\textbf{2.} Descriptive analysis of language merely reflects.... used without concern for the social prestige of these structures.
\begin{enumerate}[label=(\Alph*),leftmargin=*,itemsep=0.2em]
\item it how grammar structures and vocabulary is
\item how are grammar structures and vocabulary
\item how grammar structures and vocabulary are
\item it is how grammar structures and vocabulary are
\end{enumerate}
\par
\begingroup
\setlength{\parskip}{0pt}
\setlength{\parindent}{0pt}
\textbf{Review}\par
Answer: C\par
Confidence: high\par
Grammar rule: Embedded clause with statement word order.\par
Setelah `reflects`, dibutuhkan klausa benda dengan urutan normal, bukan inversi. Jadi bentuk yang tepat adalah `how grammar structures and vocabulary are used`.\par
\endgroup
\par
\vspace{0.25\baselineskip}
\noindent\textbf{3.} James Boswell was a Scottish author whose private journals..... personality and colorful experiences during his European journeys, meetings with renowned French philosophers, and adventures with Corsican rebels.
\begin{enumerate}[label=(\Alph*),leftmargin=*,itemsep=0.2em]
\item revelations of his fascinating
\item revealing his fascination
\item his revealing fascinations
\item reveal his fascinating
\end{enumerate}
\par
\begingroup
\setlength{\parskip}{0pt}
\setlength{\parindent}{0pt}
\textbf{Review}\par
Answer: D\par
Confidence: high\par
Grammar rule: Plural subject with simple present verb.\par
Subject `journals` berbentuk jamak, jadi perlu verb `reveal`. Pilihan (D) juga melengkapi objek dengan benar: `reveal his fascinating personality and colorful experiences`.\par
\endgroup
\par
\vspace{0.25\baselineskip}
\noindent\textbf{4.} Originally, the first European colleges consisted of groups of individuals..... joined their efforts to study sciences, medicine, and law.
\begin{enumerate}[label=(\Alph*),leftmargin=*,itemsep=0.2em]
\item who lived together and
\item whose life together and
\item and who live together
\item whose living together and
\end{enumerate}
\par
\begingroup
\setlength{\parskip}{0pt}
\setlength{\parindent}{0pt}
\textbf{Review}\par
Answer: A\par
Confidence: high\par
Grammar rule: Relative clause modifying a plural noun.\par
Frasa `groups of individuals` memerlukan relative clause yang menjelaskan orang-orang itu. (A) tepat karena membentuk `individuals who lived together and joined their efforts`.\par
\endgroup
\par
\vspace{0.25\baselineskip}
\noindent\textbf{5.} Sedimentary rock represents. geologically transformed minerals and small particles of -- --- --- deposited by the movement of water, wind, or glacial ice.
\begin{enumerate}[label=(\Alph*),leftmargin=*,itemsep=0.2em]
\item chemicals matter
\item chemical matter
\item matters of chemical
\item matters in chemicals
\end{enumerate}
\par
\begingroup
\setlength{\parskip}{0pt}
\setlength{\parindent}{0pt}
\textbf{Review}\par
Answer: B\par
Confidence: high\par
Grammar rule: Adjective + noun after preposition.\par
Setelah `particles of`, dibutuhkan noun phrase. `chemical matter` adalah susunan adjective + noun yang benar, sehingga frasa menjadi `small particles of chemical matter`.\par
\endgroup
\par
\vspace{0.25\baselineskip}
\noindent\textbf{6.} The buff pottery with red designs made by the Hohokam desert dwellers in the Southwestern regions..... for its original construction and unique shape.
\begin{enumerate}[label=(\Alph*),leftmargin=*,itemsep=0.2em]
\item are highly prized
\item is a high price
\item is highly prized
\item are prized highly
\end{enumerate}
\par
\begingroup
\setlength{\parskip}{0pt}
\setlength{\parindent}{0pt}
\textbf{Review}\par
Answer: C\par
Confidence: high\par
Grammar rule: Subject-verb agreement with singular mass noun.\par
`Pottery` adalah mass noun tunggal, jadi verb-nya harus `is`, bukan `are`. Bentuk idiomatisnya juga `is highly prized`.\par
\endgroup
\par
\vspace{0.25\baselineskip}
\noindent\textbf{7.} Soap operas and day-time comedy serials were...... originated in Chicago in the 1930s, and most consisted of 15- minute episodes broadcast five times a week.
\begin{enumerate}[label=(\Alph*),leftmargin=*,itemsep=0.2em]
\item radio shows that
\item that radio showing
\item shown that radio
\item radio that is shown
\end{enumerate}
\par
\begingroup
\setlength{\parskip}{0pt}
\setlength{\parindent}{0pt}
\textbf{Review}\par
Answer: A\par
Confidence: high\par
Grammar rule: Predicate nominative followed by relative clause.\par
Setelah `were`, kalimat memerlukan pelengkap noun phrase. (A) memberi struktur lengkap: `were radio shows that originated in Chicago`.\par
\endgroup
\par
\vspace{0.25\baselineskip}
\noindent\textbf{8.} Although food preferences vary a great deal from one individual to another, geographic location, traditional customs, and family socioeconomic position play a major role....
\begin{enumerate}[label=(\Alph*),leftmargin=*,itemsep=0.2em]
\item formational in dieta
\item informing dietary habits
\item in forming dietary habits
\item in forming dietary habitual
\end{enumerate}
\par
\begingroup
\setlength{\parskip}{0pt}
\setlength{\parindent}{0pt}
\textbf{Review}\par
Answer: C\par
Confidence: high\par
Grammar rule: Preposition + gerund after role in.\par
Pola yang benar adalah `play a role in + gerund`. Karena itu `in forming dietary habits` paling tepat secara tata bahasa.\par
\endgroup
\par
\vspace{0.25\baselineskip}
\noindent\textbf{9.} Precipitation of microscopic marine organisms on the ocean floor results in the creation of limestone, parallel or discordant layers of deposited......chemical material.
\begin{enumerate}[label=(\Alph*),leftmargin=*,itemsep=0.2em]
\item characteristics in
\item of the characteristic
\item characterized by
\item in the characterization of
\end{enumerate}
\par
\begingroup
\setlength{\parskip}{0pt}
\setlength{\parindent}{0pt}
\textbf{Review}\par
Answer: C\par
Confidence: medium\par
Grammar rule: Reduced participial phrase modifying a noun.\par
Dari pilihan yang ada, (C) paling mendekati bentuk participial phrase untuk menerangkan `layers`. Namun sumber kalimat tampak rusak sehingga hasil lengkapnya kurang mulus; item ini perlu cek manual.\par
Needs manual review: true\par
\endgroup
\par
\vspace{0.25\baselineskip}
\noindent\textbf{10.} Approximately 400 million to 250 million years ago, the ability of primitive plants to reproduce by seed......to cope with harsher habitats than could the flora that reproduced by spores.
\begin{enumerate}[label=(\Alph*),leftmargin=*,itemsep=0.2em]
\item has allowed for them
\item was allowed by them
\item have been allowed them
\item allowed them
\end{enumerate}
\par
\begingroup
\setlength{\parskip}{0pt}
\setlength{\parindent}{0pt}
\textbf{Review}\par
Answer: D\par
Confidence: high\par
Grammar rule: Simple past main verb after singular subject.\par
Subject inti kalimat adalah `the ability`, dan kalimat berlatar waktu lampau. Karena itu bentuk yang tepat adalah `allowed them to cope`.\par
\endgroup
\par
\vspace{0.25\baselineskip}
\noindent\textbf{11.} Some language experts might say......to a person speaking a language one does not understand and still determine whether the speaker is excited or exhausted, angry, or pleased.
\begin{enumerate}[label=(\Alph*),leftmargin=*,itemsep=0.2em]
\item that possibly it is to listen
\item what is possible to listen to it
\item that it is possible to listen
\item whether is it a possibility to listen
\end{enumerate}
\par
\begingroup
\setlength{\parskip}{0pt}
\setlength{\parindent}{0pt}
\textbf{Review}\par
Answer: C\par
Confidence: high\par
Grammar rule: That-clause after reporting verb with extraposed subject.\par
Setelah `might say`, dibutuhkan `that`-clause. Pola `it is possible to listen` juga benar untuk infinitive sebagai pelengkap.\par
\endgroup
\par
\vspace{0.25\baselineskip}
\noindent\textbf{12.} David Smith, who welded metal sculptures and relied on his experience as a welder in a factory during World War II, opened new paths for...... to experiment with landscape art composed of objects.
\begin{enumerate}[label=(\Alph*),leftmargin=*,itemsep=0.2em]
\item other artists
\item another, artists
\item others artists
\item the others, artist
\end{enumerate}
\par
\begingroup
\setlength{\parskip}{0pt}
\setlength{\parindent}{0pt}
\textbf{Review}\par
Answer: A\par
Confidence: high\par
Grammar rule: Object of preposition with plural noun.\par
Setelah preposisi `for`, diperlukan noun phrase yang wajar. `other artists` tepat untuk makna dan struktur: `opened new paths for other artists to experiment`.\par
\endgroup
\par
\vspace{0.25\baselineskip}
\noindent\textbf{13.} Exercise can be classified as active or passive with the former......effort and the latter the use of machines or training assistants.
\begin{enumerate}[label=(\Alph*),leftmargin=*,itemsep=0.2em]
\item involves physical
\item physics is involved
\item involving physical
\item physically involved
\end{enumerate}
\par
\begingroup
\setlength{\parskip}{0pt}
\setlength{\parindent}{0pt}
\textbf{Review}\par
Answer: C\par
Confidence: high\par
Grammar rule: With + noun + participle construction.\par
Setelah `with the former`, bentuk yang tepat adalah participle phrase untuk menjelaskan `the former`. Maka `involving physical effort` paling benar.\par
\endgroup
\par
\vspace{0.25\baselineskip}
\noindent\textbf{14.} The construction of the Alaska Highway....... in 1939 and led to a economic boom and population growth in Northern Territories.
\begin{enumerate}[label=(\Alph*),leftmargin=*,itemsep=0.2em]
\item was beginning
\item it began
\item has been begun
\item was begun
\end{enumerate}
\par
\begingroup
\setlength{\parskip}{0pt}
\setlength{\parindent}{0pt}
\textbf{Review}\par
Answer: D\par
Confidence: high\par
Grammar rule: Simple past passive.\par
Subjek `The construction of the Alaska Highway` menerima aksi, jadi diperlukan bentuk pasif. Dengan penanda waktu `in 1939`, simple past passive `was begun` adalah pilihan yang tepat.\par
\endgroup
\par
\vspace{0.25\baselineskip}
\noindent\textbf{15.} ...... the perishable quality of organic fibers, it is not possible to trace early techniques associated with weaving and interlacing of twines and threads.
\begin{enumerate}[label=(\Alph*),leftmargin=*,itemsep=0.2em]
\item Because
\item Due to the fact that
\item Because of
\item Due to the fact
\end{enumerate}
\par
\begingroup
\setlength{\parskip}{0pt}
\setlength{\parindent}{0pt}
\textbf{Review}\par
Answer: C\par
Confidence: high\par
Grammar rule: Preposition before a noun phrase of cause.\par
Bagian sesudah kosong adalah noun phrase `the perishable quality of organic fibers`, jadi penyebab harus diperkenalkan dengan preposisi `because of`, bukan `because`.\par
\endgroup
\par
\vspace{0.25\baselineskip}
\subsection*{Error Identification}
DIRECTIONS: In questions 16-40 every sentence has four words or phrases that are underlined. The four underlined portions of each sentence are marked (A), (B), (C), and (D). Identify the one word or phrase that makes the sentence incorrect. Then, on your answer sheet, find the number of the question and fill in the space that contains the letter for the answer you have chosen. Fill in the space completely.\par
\begin{exampleblock}
\textbf{EXAMPLE I}\par
Christopher Columbus \uline{has sailed}(A) from Europe in 1492 and \uline{discovered a}(B) \uline{new land}(C) he \uline{thought to}(D) be India.\par
The sentence should state, "Christopher Columbus sailed from Europe in 1492 and discovered a new land he thought to be India." Therefore, you should choose answer (A).\par
\end{exampleblock}
\begin{exampleblock}
\textbf{EXAMPLE II}\par
\uline{As the roles}(A) of people in society change, \uline{so does}(B) the \uline{rules of}(C) conduct in certain \uline{situations}(D).\par
The sentence should state, "As the roles of people in society change, so do the rules of\par
conduct in certain situations." Therefore, you should choose answer (B).\par
\end{exampleblock}
\noindent\textbf{16.} \uline{Because of}(A)the federal government \uline{did not relocate from}(B)Philadelphia to Washington, D.C., until 1800, George Washington and his wife \uline{never moved into}(C)the White House but \uline{resided in Philadelphia}(D)instead.
\par
\begingroup
\setlength{\parskip}{0pt}
\setlength{\parindent}{0pt}
\textbf{Review}\par
Answer: A\par
Confidence: high\par
Grammar rule: Because must introduce a clause; because of must be followed by a noun phrase.\par
Correct form: Because\par
Bagian (A) salah karena `because of` tidak boleh langsung diikuti klausa `the federal government did not relocate`. Bentuk yang benar adalah `Because`.\par
\endgroup
\par
\vspace{\questionblockgap}
\noindent\textbf{17.} Parallels, which are imaginary \uline{dividing lines}(A)\uline{drawing parallel}(B)to the equator and perpendicular to the Earth's axis, \uline{measure distances north}(C)and south by sectioning the globe into \uline{equally spaced circles}(D).
\par
\begingroup
\setlength{\parskip}{0pt}
\setlength{\parindent}{0pt}
\textbf{Review}\par
Answer: B\par
Confidence: high\par
Grammar rule: A reduced relative clause describing lines needs a past participle, not a present participle.\par
Correct form: drawn parallel\par
Bagian (B) salah karena frasa penjelas untuk `lines` harus berbentuk participle pasif, bukan `drawing parallel`. Bentuk yang benar adalah `drawn parallel`.\par
\endgroup
\par
\vspace{\questionblockgap}
\noindent\textbf{18.} In his novel \uline{published}(A)in 1894, Mark Twain referred to fingerprinting as a \uline{completely reliable}(B)technique \uline{for established}(C)a personal identity of twins differentiated by the \uline{shape of ridges}(D)on their fingertips.
\par
\begingroup
\setlength{\parskip}{0pt}
\setlength{\parindent}{0pt}
\textbf{Review}\par
Answer: C\par
Confidence: high\par
Grammar rule: A preposition must be followed by a gerund, not a past-tense verb form.\par
Correct form: for establishing\par
Bagian (C) salah karena setelah preposisi `for` harus digunakan gerund, bukan `established`. Bentuk yang benar adalah `for establishing`.\par
\endgroup
\par
\vspace{\questionblockgap}
\noindent\textbf{19.} Ancient Greek philosophers \uline{believed}(A)that listening to \uline{musicals arrangements}(B)and particular \uline{types of sounds}(C)led young children to develop beneficial \uline{physical and mental}(D)characteristics.
\par
\begingroup
\setlength{\parskip}{0pt}
\setlength{\parindent}{0pt}
\textbf{Review}\par
Answer: B\par
Confidence: high\par
Grammar rule: A noun used as a modifier here requires the adjective form.\par
Correct form: musical arrangements\par
Bagian (B) salah karena sebelum noun `arrangements` dibutuhkan adjective, bukan noun  plural form `musicals`. Bentuk yang benar adalah `musical arrangements`.\par
\endgroup
\par
\vspace{\questionblockgap}
\noindent\textbf{20.} At the end of the American Civil War, Andrew Jackson \uline{succeeded}(A)Abraham Lincoln \uline{as the U.S. president}(B)and \uline{after completing}(C)his term \uline{returning}(D)to his home state of Tennessee.
\par
\begingroup
\setlength{\parskip}{0pt}
\setlength{\parindent}{0pt}
\textbf{Review}\par
Answer: D\par
Confidence: high\par
Grammar rule: A coordinated predicate after and must be a finite verb parallel to the earlier verb.\par
Correct form: returned\par
Bagian (D) salah karena setelah `and after completing his term` dibutuhkan verb finite yang sejajar dengan `succeeded`, bukan participle `returning`. Bentuk yang benar adalah `returned`.\par
\endgroup
\par
\vspace{\questionblockgap}
\noindent\textbf{21.} Before the American Revolution, the task of \uline{maintaining law and order}(A)\uline{was carrying out}(B)by \uline{elected officials}(C)and citizens to keep a careful watch for fires and suspicious activities in towns and \uline{rural areas alike}(D).
\par
\begingroup
\setlength{\parskip}{0pt}
\setlength{\parindent}{0pt}
\textbf{Review}\par
Answer: B\par
Confidence: high\par
Grammar rule: Passive voice requires be plus past participle.\par
Correct form: was carried out\par
Bagian (B) salah karena bentuk pasif harus memakai past participle setelah `was`, bukan `carrying out`. Bentuk yang benar adalah `was carried out`.\par
\endgroup
\par
\vspace{\questionblockgap}
\noindent\textbf{22.} Edmund Malone was a great scholar \uline{whom depicted}(A)the life and work of William Shakespeare and \uline{contributed his expertise}(B)in language \uline{to publishing a formal}(C)biography of the playwright and his \uline{humble beginnings}(D).
\par
\begingroup
\setlength{\parskip}{0pt}
\setlength{\parindent}{0pt}
\textbf{Review}\par
Answer: A\par
Confidence: high\par
Grammar rule: A subject relative pronoun must be who, not whom.\par
Correct form: who depicted\par
Bagian (A) salah karena pronoun relatif yang menjadi subjek verb `depicted` harus `who`, bukan `whom`. Bentuk yang benar adalah `who depicted`.\par
\endgroup
\par
\vspace{\questionblockgap}
\noindent\textbf{23.} Water levels in the Great Lakes \uline{fluctuate dramatically}(A)over periods of years by \uline{as much as one}(B)meter, and during the spring runoff, the levels \uline{may rise}(C)or fall by as much as two meters, particularly in Lake Erie, \uline{the shallow of the five}(D).
\par
\begingroup
\setlength{\parskip}{0pt}
\setlength{\parindent}{0pt}
\textbf{Review}\par
Answer: D\par
Confidence: high\par
Grammar rule: A superlative adjective is required when comparing one item with all five.\par
Correct form: the shallowest of the five\par
Bagian (D) salah karena perbandingan dengan seluruh kelompok lima danau memerlukan bentuk superlatif, bukan `shallow`. Bentuk yang benar adalah `the shallowest of the five`.\par
\endgroup
\par
\vspace{\questionblockgap}
\noindent\textbf{24.} Whether \uline{sounds are perceived}(A)to be musical or not depends on \uline{learning more}(B)and \uline{cultural conditioning than}(C)the physical properties of sounds or \uline{innate characteristics}(D)of the human brain.
\par
\begingroup
\setlength{\parskip}{0pt}
\setlength{\parindent}{0pt}
\textbf{Review}\par
Answer: B\par
Confidence: high\par
Grammar rule: The comparative adverb more must modify depends on, not the noun phrase after it.\par
Correct form: more on learning\par
Bagian (B) salah karena `more` harus menerangkan hubungan `depends on`, bukan diletakkan sesudah `learning`. Bentuk yang benar adalah `more on learning`.\par
\endgroup
\par
\vspace{\questionblockgap}
\noindent\textbf{25.} Baobab trees \uline{that closely related}(A)to the family of bottle and cactus trees have thick trunks with soft centers \uline{that store water}(B)in \uline{sufficient amounts}(C)to enable the plants \uline{to weather long}(D)periods of desert droughts.
\par
\begingroup
\setlength{\parskip}{0pt}
\setlength{\parindent}{0pt}
\textbf{Review}\par
Answer: A\par
Confidence: high\par
Grammar rule: A relative clause needs a form of be before the adjective related.\par
Correct form: are closely related\par
Bagian (A) salah karena klausa relatif memerlukan linking verb sebelum adjective `related`. Bentuk yang benar adalah `are closely related`.\par
\endgroup
\par
\vspace{\questionblockgap}
\noindent\textbf{26.} Since \uline{the early twentieth century}(A), the mass-production of automobiles \uline{has begun to change}(B)the American lifestyle and landscape, \uline{as ready access by}(C)a motor vehicle determines \uline{where do people}(D)build houses and other facilities.
\par
\begingroup
\setlength{\parskip}{0pt}
\setlength{\parindent}{0pt}
\textbf{Review}\par
Answer: D\par
Confidence: high\par
Grammar rule: An indirect question uses statement word order, not auxiliary inversion.\par
Correct form: where people\par
Bagian (D) salah karena indirect question setelah `determines` tidak memakai inversi `do people`. Bentuk yang benar adalah `where people`.\par
\endgroup
\par
\vspace{\questionblockgap}
\noindent\textbf{27.} Although radiocarbon dating of \uline{archeological finds}(A)can prove to be \uline{fair reliable}(B)for establishing the age of prehistoric objects, additional techniques must \uline{be employed}(C)to date objects \uline{older than}(D)50,000 years.
\par
\begingroup
\setlength{\parskip}{0pt}
\setlength{\parindent}{0pt}
\textbf{Review}\par
Answer: B\par
Confidence: high\par
Grammar rule: An adjective must be modified by an adverb, not another adjective.\par
Correct form: fairly reliable\par
Bagian (B) salah karena adjective `reliable` harus diterangkan oleh adverb, bukan adjective `fair`. Bentuk yang benar adalah `fairly reliable`.\par
\endgroup
\par
\vspace{\questionblockgap}
\noindent\textbf{28.} \uline{Owing}(A)to their \uline{superior skill}(B), highly competitive athletes \uline{have been known}(C)to win contests and break records even \uline{when suffered}(D)from injuries, physical disorders, and infections.
\par
\begingroup
\setlength{\parskip}{0pt}
\setlength{\parindent}{0pt}
\textbf{Review}\par
Answer: D\par
Confidence: high\par
Grammar rule: After when in a reduced clause, the active meaning requires a present participle.\par
Correct form: when suffering\par
Bagian (D) salah karena setelah `when` pada klausa tereduksi aktif diperlukan present participle, bukan past form `suffered`. Bentuk yang benar adalah `when suffering`.\par
\endgroup
\par
\vspace{\questionblockgap}
\noindent\textbf{29.} The keyboard \uline{was invented}(A)in medieval Europe during the \uline{rise of}(B)multipart musical compositions and \uline{replacement}(C)the mechanical dulcimer in \uline{domestic uses}(D)and public performances.
\par
\begingroup
\setlength{\parskip}{0pt}
\setlength{\parindent}{0pt}
\textbf{Review}\par
Answer: C\par
Confidence: high\par
Grammar rule: Parallel noun phrases after and require the noun replacement to take its complement with of.\par
Correct form: the replacement of\par
Bagian (C) salah karena noun `replacement` memerlukan pelengkap dengan preposisi `of` agar sejajar dengan `the rise of`. Bentuk yang benar adalah `the replacement of`.\par
\endgroup
\par
\vspace{\questionblockgap}
\noindent\textbf{30.} The Paleo-Indian tribes \uline{were able to switch}(A)from food gathering and nomadic way of life to permanent settlements \uline{only they}(B)acquired \uline{means of planting}(C)and harvesting \uline{more predictable}(D)crops of grain and beans.
\par
\begingroup
\setlength{\parskip}{0pt}
\setlength{\parindent}{0pt}
\textbf{Review}\par
Answer: B\par
Confidence: high\par
Grammar rule: A subordinating conjunction is needed before the clause after only.\par
Correct form: only after they\par
Bagian (B) salah karena setelah `only` dibutuhkan konjungsi subordinatif untuk menghubungkan klausa waktu. Bentuk yang benar adalah `only after they`.\par
\endgroup
\par
\vspace{\questionblockgap}
\noindent\textbf{31.} Hand axes \uline{founded}(A)in the high banks of the Thames River \uline{have been identified}(B)as much older artifacts than pottery shards \uline{excavated in}(C)the lower flood planes \uline{located in different}(D)geological strata.
\par
\begingroup
\setlength{\parskip}{0pt}
\setlength{\parindent}{0pt}
\textbf{Review}\par
Answer: A\par
Confidence: high\par
Grammar rule: A reduced passive clause requires the past participle found, not founded.\par
Correct form: found\par
Bagian (A) salah karena frasa tereduksi pasif untuk `Hand axes` harus memakai past participle `found`, bukan `founded`. Bentuk yang benar adalah `found`.\par
\endgroup
\par
\vspace{\questionblockgap}
\noindent\textbf{32.} In \uline{stained-glass work}(A), cut pieces of colored glass \uline{are positioned}(B)in incremental frames made of \uline{shape metal strips}(C)and soldered \uline{to a large overall}(D)frame.
\par
\begingroup
\setlength{\parskip}{0pt}
\setlength{\parindent}{0pt}
\textbf{Review}\par
Answer: C\par
Confidence: high\par
Grammar rule: A past participle adjective is needed to modify metal strips.\par
Correct form: shaped metal strips\par
Bagian (C) salah karena sebelum noun `metal strips` dibutuhkan adjective `shaped`, bukan noun atau verb form `shape`. Bentuk yang benar adalah `shaped metal strips`.\par
\endgroup
\par
\vspace{\questionblockgap}
\noindent\textbf{33.} The colonies of settlers \uline{who were called}(A)the "Pennsylvania Dutch" actually \uline{consisted from}(B)German, Swiss, and Moravian members of religious sects \uline{who secluded}(C)\uline{themselves from}(D)political and social affairs of their secular neighbors.
\par
\begingroup
\setlength{\parskip}{0pt}
\setlength{\parindent}{0pt}
\textbf{Review}\par
Answer: B\par
Confidence: high\par
Grammar rule: The verb consist takes the preposition of, not from.\par
Correct form: consisted of\par
Bagian (B) salah karena verb `consist` harus diikuti preposisi `of`, bukan `from`. Bentuk yang benar adalah `consisted of`.\par
\endgroup
\par
\vspace{\questionblockgap}
\noindent\textbf{34.} The price theory \uline{that represents}(A)the core of microeconomics \uline{explained}(B)how the variability of supply and demand \uline{in competitive}(C)markets \uline{creates}(D)an interplay of prices for goods and services.
\par
\begingroup
\setlength{\parskip}{0pt}
\setlength{\parindent}{0pt}
\textbf{Review}\par
Answer: B\par
Confidence: high\par
Grammar rule: A general present statement needs the present-tense verb to match the surrounding clause.\par
Correct form: explains\par
Bagian (B) salah karena kalimat ini menyatakan fakta umum dan harus sejajar dengan present tense `represents` dan `creates`. Bentuk yang benar adalah `explains`.\par
\endgroup
\par
\vspace{\questionblockgap}
\noindent\textbf{35.} Historians believe that Navajo tribes \uline{migrated to}(A)the south from the northwestern \uline{coasts of}(B)Alaska and Canada \uline{and populated the}(C)Southwestern states \uline{as recent as}(D)500 years ago.
\par
\begingroup
\setlength{\parskip}{0pt}
\setlength{\parindent}{0pt}
\textbf{Review}\par
Answer: D\par
Confidence: high\par
Grammar rule: The idiomatic time expression is as recently as, not as recent as.\par
Correct form: as recently as\par
Bagian (D) salah karena ungkapan waktu yang benar adalah `as recently as`, bukan `as recent as`. Bentuk yang benar adalah `as recently as`.\par
\endgroup
\par
\vspace{\questionblockgap}
\noindent\textbf{36.} During his remarkable career, Paul Samuelson \uline{served as}(A)a magazine columnist, \uline{consultant to}(B)research organizations, \uline{advised two}(C)presidents, and lecturer on mathematical analysis \uline{on levels of}(D)employment.
\par
\begingroup
\setlength{\parskip}{0pt}
\setlength{\parindent}{0pt}
\textbf{Review}\par
Answer: C\par
Confidence: high\par
Grammar rule: Items in a noun list after served as should be parallel noun phrases, not a finite verb.\par
Correct form: adviser to two\par
Bagian (C) salah karena unsur deret setelah `served as` harus paralel dalam bentuk noun phrase, bukan verb `advised`. Bentuk yang benar adalah `adviser to two`.\par
\endgroup
\par
\vspace{\questionblockgap}
\noindent\textbf{37.} In 1838, Auguste Compte devised the term "sociology" \uline{to refer}(A)to a new science that \uline{will discover}(B)regularities of human society \uline{similar to}(C)laws of nature \uline{by employing}(D)investigative methods.
\par
\begingroup
\setlength{\parskip}{0pt}
\setlength{\parindent}{0pt}
\textbf{Review}\par
Answer: B\par
Confidence: high\par
Grammar rule: A future action viewed from a past reporting point takes would, not will.\par
Correct form: would discover\par
Bagian (B) salah karena dalam pelaporan dari sudut waktu lampau `devised`, future-in-the-past lebih tepat memakai `would`, bukan `will`. Bentuk yang benar adalah `would discover`.\par
\endgroup
\par
\vspace{\questionblockgap}
\noindent\textbf{38.} Although zebras are \uline{related to horses}(A), \uline{they are size}(B)is closer to that of donkeys \uline{whom they}(C)resemble in behavior and body form, \uline{with short legs}(D), high shoulders, a stiff black mane, large ears, and a tufted tail.
\par
\begingroup
\setlength{\parskip}{0pt}
\setlength{\parindent}{0pt}
\textbf{Review}\par
Answer: B\par
Confidence: high\par
Grammar rule: A possessive determiner is needed before size.\par
Correct form: their size\par
Bagian (B) salah karena sebelum noun `size` dibutuhkan possessive determiner, bukan struktur `they are size`. Bentuk yang benar adalah `their size`.\par
\endgroup
\par
\vspace{\questionblockgap}
\noindent\textbf{39.} \uline{Influenced}(A)by the advancements in behavioral sciences, George Mead \uline{argued}(B)that communication and interaction \uline{with the environment}(C)are the keys \uline{to understand}(D)an organism's self-consciousness.
\par
\begingroup
\setlength{\parskip}{0pt}
\setlength{\parindent}{0pt}
\textbf{Review}\par
Answer: D\par
Confidence: high\par
Grammar rule: After the noun keys, to functions as a preposition and should be followed by a gerund.\par
Correct form: to understanding\par
Bagian (D) salah karena pada frasa `the keys to ...`, `to` adalah preposisi sehingga harus diikuti gerund, bukan infinitive `to understand`. Bentuk yang benar adalah `to understanding`.\par
\endgroup
\par
\vspace{\questionblockgap}
\noindent\textbf{40.} Among the sixteenth-century Dutch painters, Peter Bruegel is \uline{the most}(A)notable \uline{because of his}(B)portrayal of peasant scenes and proverbial folktales \uline{that have charmed}(C)art aficionados \uline{since more}(D)than 400 years.
\par
\begingroup
\setlength{\parskip}{0pt}
\setlength{\parindent}{0pt}
\textbf{Review}\par
Answer: D\par
Confidence: high\par
Grammar rule: A duration with present perfect uses for, not since, when no starting point is given.\par
Correct form: for more than 400 years\par
Bagian (D) salah karena present perfect untuk durasi membutuhkan `for` jika yang disebut adalah lamanya waktu, bukan titik awal. Bentuk yang benar adalah `for more than 400 years`.\par
\endgroup
\par
\vspace{\questionblockgap}
\section*{Reading Comprehension}
\begin{center}
\textbf{Time 55 minutes}
\end{center}
DIRECTIONS: In this section you will read several passages. Each passage is followed by a series of questions. For questions 1-50, you need to select the best answer, (A), (B), (C), or (D), to each question. Then, on your answer sheet, find the number of the question and fill in the space that contains the letter of the answer you have selected. Fill in the space completely. Answer all questions following a passage on the basis of what is stated or implied in the passage.\par
\begin{exampleblock}
\small
\sloppy
Read the following passage:\par
A tomahawk is a small ax used as a tool and a weapon by the North American Indian\par
tribes. An average tomahawk was not very long and did not weigh a great deal. Originally,\par
the head of the tomahawk was made of a shaped stone or an animal bone and was mounted on\par
Line a wooden handle. After the arrival of the European settlers, the Indians began to use toma-\par
hawks with iron heads. Indian males and females of all ages used tomahawks to chop and cut wood, pound stakes into the ground to put up wigwams, and do many other chores. Indian warriors relied on tomahawks as weapons and even threw them at their enemies. Some types of tomahawks were used in religious ceremonies. Contemporary American idioms reflect\par
this aspect of American heritage.\par
\end{exampleblock}
\begin{exampleblock}
\textbf{EXAMPLE I}\par
\noindent Early tomahawk heads were made of
\begin{enumerate}[label=(\Alph*),leftmargin=*,itemsep=0.2em]
\item stone or bone
\item wood or sticks
\item European iron
\item religious weapons
\end{enumerate}
\vspace{0.25\baselineskip}
According to the passage, early tomahawk heads were made of stone or bone. Therefore, the correct answer is (A).\par
\end{exampleblock}
\begin{exampleblock}
\textbf{EXAMPLE II}\par
\noindent How has the Indian use of tomahawks affected American daily life today?
\begin{enumerate}[label=(\Alph*),leftmargin=*,itemsep=0.2em]
\item Tomahawks are still used as weapons.
\item Tomahawks are used as tools for.certain jobs.
\item Contemporary language refers to tomahawks.
\item Indian tribes cherish tomahawks as heirlooms.
\end{enumerate}
\vspace{0.25\baselineskip}
The passage states, "Contemporary American idioms reflect this aspect of American heritage." The correct answer is (C). After you read the directions, begin work on the questions.\par
\end{exampleblock}
\subsection*{Questions 1-12}
\begingroup
\small
\sloppy
\setlength{\parskip}{0pt}
\setlength{\parindent}{0pt}
Architecture has social purposes and meets practical needs by means of combining art\par
and technological innovations. In building construction, however, an emergence of new\par
materials does not make its precursors obsolete, and architectural knowledge is cumulative.\par
The fact that today much is constructed from prefabricated concrete does not do away with\par
brick. Furthermore, despite dramatic changes and increased technological sophistication of\par
architectural design and construction, the essential apparatus of erecting a building has\par
remained rooted in preindustrial traditional practices passed down during the millennia. The\par
social and utilitarian expectations of structures are largely based on elemental demands of\par
keeping out elements and enemies, ameliorating the extremes of heat, and avoiding the\par
intrusion of wind, precipitation, and pests.\par
Gravity, air pressure, and earthquakes can induce tensions that have to be accounted for\par
when constructing functional enclosed space. Vertical stacking of masonry materials causes\par
compression that can lead to important problems when a structure is spanned to build a\par
roof and connect walls. Arches, vaults, and domes were specifically developed to alleviate\par
the compression by directing the spanning element along a curve rather than a straight line.\par
Building suspension structures, dams, and tunnels became possible in the nineteenth\par
century with the increased availability of steel that could reinforce structural frames and\par
enable. them to withstand natural forces previously believed to be insurmountable.\par
Functional evolutions of modern buildings create new demands on the analysis of struc-\par
tural behavior and engineering. Few occupants of skyscrapers view elevators as elaborate\par
systems of vertical transportation. Humidity and temperature control, forced ventilation, nat-\par
ural and artificial lighting, sanitation and disposal of waste, electrical wiring, and fire pre-\par
vention make very tall constructions engineering marvels that also must be aesthetically\par
pleasing and physically convenient.\par
Erecting a structure involves a great deal more than merely attending to the aesthetics and\par
psychological experience of architectural space. The shape, size, and incombustibility of locally\par
available construction materials fostered developments of specific technologies, and brick and\par
stone masonry have evolved in. response to the need for structural durability. Advances in civil\par
engineering and knowledge associated with properties of building materials combine to lead to\par
innovations in architectural design. Tools and skills required to exploit easily obtainable materi-\par
als have continued to inform the development of modern industrialized technologies.\par
\endgroup
\medskip
\noindent\textbf{1.} What is the main topic of the passage?
\begin{enumerate}[label=(\Alph*),leftmargin=*,itemsep=0.2em]
\item The modern art of architecture and social pressure
\item The profound importance of tradition in architecture
\item The mutual impact of architecture and technology
\item The great technological advances in building materials
\end{enumerate}
\par
\begingroup
\setlength{\parskip}{0pt}
\setlength{\parindent}{0pt}
\textbf{Review}\par
Answer: C\par
Confidence: high\par
Evidence refs: lines 1-2, lines 29-31\par
Evidence quote: "Architecture has social purposes and meets practical needs by means of combining art and technological innovations... Advances in civil engineering and knowledge associated with properties of building materials combine to lead to innovations in architectural design."\par
Jawaban (C) karena dari awal sampai akhir passage membahas hubungan timbal balik antara arsitektur dan teknologi: arsitektur memakai inovasi teknologi, sementara perkembangan material dan engineering mendorong desain arsitektur baru.\par
\endgroup
\par
\vspace{0.25\baselineskip}
\noindent\textbf{2.} The word "obsolete" in line 3 is closest in meaning to
\begin{enumerate}[label=(\Alph*),leftmargin=*,itemsep=0.2em]
\item obvious
\item obstinate
\item antiquarian
\item antiquated
\end{enumerate}
\par
\begingroup
\setlength{\parskip}{0pt}
\setlength{\parindent}{0pt}
\textbf{Review}\par
Answer: D\par
Confidence: high\par
Evidence refs: lines 2-4\par
Evidence quote: "an emergence of new materials does not make its precursors obsolete"\par
Jawaban (D) karena dalam konteks ini obsolete berarti tidak lagi dipakai atau ketinggalan zaman, paling dekat dengan antiquated.\par
\endgroup
\par
\vspace{0.25\baselineskip}
\noindent\textbf{3.} The author mentions the word "brick" in line 5 as an example of which of the following?
\begin{enumerate}[label=(\Alph*),leftmargin=*,itemsep=0.2em]
\item How old techniques can continue to remain practical
\item How old buildings can coexist with modern architecture
\item How new knowledge can supplant traditional technology
\item How new design can improve traditional construction
\end{enumerate}
\par
\begingroup
\setlength{\parskip}{0pt}
\setlength{\parindent}{0pt}
\textbf{Review}\par
Answer: A\par
Confidence: high\par
Evidence refs: lines 3-5\par
Evidence quote: "new materials does not make its precursors obsolete... much is constructed from prefabricated concrete does not do away with brick"\par
Jawaban (A) karena brick dipakai sebagai contoh bahwa teknik atau material lama tetap praktis meskipun material baru sudah muncul.\par
\endgroup
\par
\vspace{0.25\baselineskip}
\noindent\textbf{4.} It can be inferred from the passage that pragmatic requirements of buildings
\begin{enumerate}[label=(\Alph*),leftmargin=*,itemsep=0.2em]
\item retain essential sophistication
\item hold constant over time
\item stagnate over millennia
\item stay rooted in the elements
\end{enumerate}
\par
\begingroup
\setlength{\parskip}{0pt}
\setlength{\parindent}{0pt}
\textbf{Review}\par
Answer: B\par
Confidence: high\par
Evidence refs: lines 7-10\par
Evidence quote: "The social and utilitarian expectations of structures are largely based on elemental demands of keeping out elements and enemies..."\par
Jawaban (B) karena kebutuhan praktis bangunan digambarkan tetap bertumpu pada tuntutan dasar yang sama, seperti perlindungan dari cuaca dan ancaman, sehingga relatif konstan sepanjang waktu.\par
\endgroup
\par
\vspace{0.25\baselineskip}
\noindent\textbf{5.} What can be inferred from the passage about reducing the effects of material compacting?
\begin{enumerate}[label=(\Alph*),leftmargin=*,itemsep=0.2em]
\item Masonry is stacked vertically by increasing compression.
\item Downward pressure is dispersed by semicircular roofs.
\item Buildings are spanned to account for the force of gravity.
\item Vertical roofs are bent to counteract air pressure.
\end{enumerate}
\par
\begingroup
\setlength{\parskip}{0pt}
\setlength{\parindent}{0pt}
\textbf{Review}\par
Answer: B\par
Confidence: high\par
Evidence refs: lines 12-15\par
Evidence quote: "Vertical stacking of masonry materials causes compression... Arches, vaults, and domes were specifically developed to alleviate the compression by directing the spanning element along a curve"\par
Jawaban (B) karena passage menjelaskan bahwa tekanan kompresi dikurangi dengan bentuk lengkung seperti arches, vaults, dan domes, jadi gaya ke bawah disebarkan melalui struktur melengkung.\par
\endgroup
\par
\vspace{0.25\baselineskip}
\noindent\textbf{6.} The word "withstand" in line 1 8 is closest in meaning to
\begin{enumerate}[label=(\Alph*),leftmargin=*,itemsep=0.2em]
\item endure
\item enlarge
\item withdraw
\item withhold
\end{enumerate}
\par
\begingroup
\setlength{\parskip}{0pt}
\setlength{\parindent}{0pt}
\textbf{Review}\par
Answer: A\par
Confidence: high\par
Evidence refs: lines 16-18\par
Evidence quote: "steel that could reinforce structural frames and enable them to withstand natural forces previously believed to be insurmountable"\par
Jawaban (A) karena withstand natural forces berarti mampu menahan atau endure kekuatan alam.\par
\endgroup
\par
\vspace{0.25\baselineskip}
\noindent\textbf{7.} The purpose of paragraph 3 is to suggest that
\begin{enumerate}[label=(\Alph*),leftmargin=*,itemsep=0.2em]
\item tall buildings require large amounts of wiring to make them functional
\item architectural innovations pose new challenges for technological development
\item skyscrapers need to be appealing and convenient for their occupants
\item architects of modern buildings create a demand for engineering talent
\end{enumerate}
\par
\begingroup
\setlength{\parskip}{0pt}
\setlength{\parindent}{0pt}
\textbf{Review}\par
Answer: B\par
Confidence: high\par
Evidence refs: lines 19-24\par
Evidence quote: "Functional evolutions of modern buildings create new demands on the analysis of structural behavior and engineering... make very tall constructions engineering marvels"\par
Jawaban (B) karena bagian ini menekankan bahwa perkembangan fungsi bangunan modern menimbulkan tuntutan teknik baru, jadi inovasi arsitektur memunculkan tantangan baru bagi pengembangan teknologi dan engineering.\par
\endgroup
\par
\vspace{0.25\baselineskip}
\noindent\textbf{8.} The word "marvels" in line 23 is closest in meaning to
\begin{enumerate}[label=(\Alph*),leftmargin=*,itemsep=0.2em]
\item miracle
\item mirage
\item conception
\item construction
\end{enumerate}
\par
\begingroup
\setlength{\parskip}{0pt}
\setlength{\parindent}{0pt}
\textbf{Review}\par
Answer: A\par
Confidence: high\par
Evidence refs: lines 21-23\par
Evidence quote: "Humidity and temperature control... fire prevention make very tall constructions engineering marvels"\par
Jawaban (A) karena marvels dalam frasa engineering marvels berarti sesuatu yang mengagumkan atau wonders, paling dekat dengan miracle.\par
\endgroup
\par
\vspace{0.25\baselineskip}
\noindent\textbf{9.} According to the passage, what is one of the important requirements of building materials?
\begin{enumerate}[label=(\Alph*),leftmargin=*,itemsep=0.2em]
\item They need to be large and well shaped.
\item They should be locally produced.
\item They do not affect human psychology.
\item They have to be inflammable.
\end{enumerate}
\par
\begingroup
\setlength{\parskip}{0pt}
\setlength{\parindent}{0pt}
\textbf{Review}\par
Answer: B\par
Confidence: medium\par
Evidence refs: lines 26-28\par
Evidence quote: "The shape, size, and incombustibility of locally available construction materials fostered developments of specific technologies"\par
Jawaban (B) karena passage secara eksplisit menyebut locally available construction materials sebagai faktor penting. Pilihan lain tidak cocok: inflammable berlawanan dengan incombustibility, dan large tidak disebut sebagai keharusan.\par
\endgroup
\par
\vspace{0.25\baselineskip}
\noindent\textbf{10.} The word "fostered" in line 27 is closest in meaning to
\begin{enumerate}[label=(\Alph*),leftmargin=*,itemsep=0.2em]
\item founded
\item focused
\item encouraged
\item enveloped
\end{enumerate}
\par
\begingroup
\setlength{\parskip}{0pt}
\setlength{\parindent}{0pt}
\textbf{Review}\par
Answer: C\par
Confidence: high\par
Evidence refs: lines 26-28\par
Evidence quote: "The shape, size, and incombustibility of locally available construction materials fostered developments of specific technologies"\par
Jawaban (C) karena fostered developments berarti mendorong atau encouraged perkembangan teknologi tertentu.\par
\endgroup
\par
\vspace{0.25\baselineskip}
\noindent\textbf{11.} It can be inferred from the passage that architecture and engineering
\begin{enumerate}[label=(\Alph*),leftmargin=*,itemsep=0.2em]
\item are at the opposite ends of the technological spectrum
\item go hand in hand to promote art and science
\item compete for technological advancements
\item supercede aesthetic and experiential values
\end{enumerate}
\par
\begingroup
\setlength{\parskip}{0pt}
\setlength{\parindent}{0pt}
\textbf{Review}\par
Answer: B\par
Confidence: high\par
Evidence refs: lines 1-2, lines 28-30\par
Evidence quote: "Architecture has social purposes and meets practical needs by means of combining art and technological innovations... Advances in civil engineering... combine to lead to innovations in architectural design."\par
Jawaban (B) karena passage menunjukkan arsitektur dan engineering bekerja saling mendukung: arsitektur memadukan seni dengan teknologi, dan engineering membantu melahirkan inovasi desain.\par
\endgroup
\par
\vspace{0.25\baselineskip}
\noindent\textbf{12.} The word "inform" in line 31 is closest in meaning to
\begin{enumerate}[label=(\Alph*),leftmargin=*,itemsep=0.2em]
\item insist on
\item infringe on
\item contrast with
\item contribute to
\end{enumerate}
\par
\begingroup
\setlength{\parskip}{0pt}
\setlength{\parindent}{0pt}
\textbf{Review}\par
Answer: D\par
Confidence: high\par
Evidence refs: lines 30-31\par
Evidence quote: "Tools and skills required to exploit easily obtainable materials have continued to inform the development of modern industrialized technologies."\par
Jawaban (D) karena inform the development di sini berarti ikut membentuk atau contribute to perkembangan teknologi modern.\par
\endgroup
\par
\vspace{0.25\baselineskip}
\subsection*{Questions 13-21}
\begingroup
\small
\sloppy
\setlength{\parskip}{0pt}
\setlength{\parindent}{0pt}
In the winter of 1848, James Marshall, a carpenter from New Jersey, was contracted by\par
John Sutter to build a sawmill in Coloma in eastern California. On January 24, Marshall dis-\par
covered gold on the mill grounds. He took his find to Sutter, who attempted to keep the\par
news a secret and made his hired hands promise not to divulge it. However, Samuel\par
Brannan, a shrewd San Francisco businessman, who owned a general store in the vicinity\par
of Sutter's mill, saw potential for profit. In the spring, after hearing the rumors from the\par
storekeeper in Coloma, Brannan hastily traveled to the Sacramento Valley and on his return\par
widely disseminated the news about the untold riches in the diggings.\par
The first new arrivals assumed that the story of the gold was a hoax, and many were idle\par
curiosity seekers. However, within weeks, scores of city dwellers began to move by ships\par
and boats to reach the Coloma fort, and in subsequent months, rumors spread further down\par
the coast of California, across the Mexican border, and to Hawaii. In eastern states, the\par
press published reports of miners who struck it rich overnight so vast were the California\par
gold deposits. Between 1848 and 1851, in all, around half a million fortune-seekers from\par
around the globe descended o n California to pursue instant wealth and thus took part in\par
one of the largest migrations in history.\par
The accounts of the Sutter Mill deposits were not entirely untruthful. In December of\par
1848, President James Polk confirmed the news about the California gold in a statement to\par
Congress, and the message was picked up by press agencies worldwide. Not all Forty-\par
Niners, as the prospectors came to be known, made their journey in vain. The swath of\par
gold-bearing quartz ran 100 miles wide along the Sierra Nevada mountain range and\par
stretched from La Porte to Mariposa. During the initial stages of the Gold Fever, yields of\par
\$300 to \$400 a day were not infrequent, and b y the end of 1849, nearly \$10 million worth\par
of deposits had been excavated. However, a s competition for land claims increased, yields\par
and profits began to drop off, and ore nuggets and dust began to disappear from dry\par
streambeds and ravines. As mining operations became more laborious, individual prospec-\par
tors were replaced by large commercial ventures, when groups of miners employed their\par
penniless counterparts to divert and dry up rivers with strategically positioned dams and\par
canals. The gold rush of 1849 was followed by the next wave of migrants headed to the\par
Klondike region in Yukon Territory in mid-1897 to recover the gold at the conflation of two\par
rivers\par
\endgroup
\medskip
\noindent\textbf{13.} The word "contracted" in line 1 is closest in meaning to
\begin{enumerate}[label=(\Alph*),leftmargin=*,itemsep=0.2em]
\item empowered
\item employed
\item emboldened
\item embodied
\end{enumerate}
\par
\begingroup
\setlength{\parskip}{0pt}
\setlength{\parindent}{0pt}
\textbf{Review}\par
Answer: B\par
Confidence: high\par
Evidence refs: lines 1-2\par
Evidence quote: "James Marshall ... was contracted by John Sutter to build a sawmill"\par
Jawaban (B) karena dalam konteks ini contracted by berarti dipekerjakan untuk melakukan pekerjaan tertentu, bukan empowered atau embodied.\par
\endgroup
\par
\vspace{0.25\baselineskip}
\noindent\textbf{14.} According to the passage, who was initially instrumental in diffusing accounts of Coloma gold?
\begin{enumerate}[label=(\Alph*),leftmargin=*,itemsep=0.2em]
\item A carpenter
\item A mill owner
\item A merchant
\item A business operator
\end{enumerate}
\par
\begingroup
\setlength{\parskip}{0pt}
\setlength{\parindent}{0pt}
\textbf{Review}\par
Answer: C\par
Confidence: high\par
Evidence refs: lines 5-8\par
Evidence quote: "Samuel Brannan ... owned a general store ... widely disseminated the news about the untold riches"\par
Jawaban (C) karena orang yang mula-mula menyebarluaskan berita emas adalah Samuel Brannan, dan passage mengenalkannya sebagai pemilik general store, yaitu seorang merchant.\par
\endgroup
\par
\vspace{0.25\baselineskip}
\noindent\textbf{15.} In line 5, the word "vicinity" is closest in meaning to
\begin{enumerate}[label=(\Alph*),leftmargin=*,itemsep=0.2em]
\item proximity
\item propensity
\item propriety
\item profundity
\end{enumerate}
\par
\begingroup
\setlength{\parskip}{0pt}
\setlength{\parindent}{0pt}
\textbf{Review}\par
Answer: A\par
Confidence: high\par
Evidence refs: line 5\par
Evidence quote: "a general store in the vicinity of Sutter's mill"\par
Jawaban (A) karena vicinity berarti area yang dekat atau nearby area, sehingga paling dekat maknanya dengan proximity.\par
\endgroup
\par
\vspace{0.25\baselineskip}
\noindent\textbf{16.} The purpose of paragraph 2 is to demonstrate
\begin{enumerate}[label=(\Alph*),leftmargin=*,itemsep=0.2em]
\item the suspicions of curiosity seekers
\item the significance of population resettlement
\item the influence of published press reports
\item the wealth of gold deposits
\end{enumerate}
\par
\begingroup
\setlength{\parskip}{0pt}
\setlength{\parindent}{0pt}
\textbf{Review}\par
Answer: B\par
Confidence: high\par
Evidence refs: lines 10-16\par
Evidence quote: "scores of city dwellers began to move ... around half a million fortune-seekers ... took part in one of the largest migrations in history"\par
Jawaban (B) karena bagian ini menekankan perpindahan penduduk dalam skala sangat besar, dari puluhan orang sampai setengah juta migran, jadi fungsinya menunjukkan signifikansi resettlement itu.\par
\endgroup
\par
\vspace{0.25\baselineskip}
\noindent\textbf{17.} The word "scores" in line 10 is closest in meaning to
\begin{enumerate}[label=(\Alph*),leftmargin=*,itemsep=0.2em]
\item a greater good
\item a great stock
\item a great many
\item a greater evil
\end{enumerate}
\par
\begingroup
\setlength{\parskip}{0pt}
\setlength{\parindent}{0pt}
\textbf{Review}\par
Answer: C\par
Confidence: high\par
Evidence refs: line 10\par
Evidence quote: "within weeks, scores of city dwellers began to move"\par
Jawaban (C) karena scores of dalam konteks jumlah orang berarti many atau a great many.\par
\endgroup
\par
\vspace{0.25\baselineskip}
\noindent\textbf{18.} It can be inferred from the passage that California gold had a variety of
\begin{enumerate}[label=(\Alph*),leftmargin=*,itemsep=0.2em]
\item polemical and ecological outcomes
\item proclamations and denunciations in the White House
\item proponents and opponents in Congress
\item political and economic ramifications
\end{enumerate}
\par
\begingroup
\setlength{\parskip}{0pt}
\setlength{\parindent}{0pt}
\textbf{Review}\par
Answer: D\par
Confidence: medium\par
Evidence refs: lines 14-16, lines 18-24\par
Evidence quote: "around half a million fortune-seekers ... one of the largest migrations in history ... President James Polk confirmed the news ... yields of \$300 to \$400 a day"\par
Jawaban (D) karena passage menunjukkan dampak politik saat Presiden Polk dan Congress terlibat dalam konfirmasi berita, serta dampak ekonomi lewat arus pencari kekayaan, hasil tambang, dan penurunan profit.\par
\endgroup
\par
\vspace{0.25\baselineskip}
\noindent\textbf{19.} In line 20, the word "prospectors" is closest in meaning to
\begin{enumerate}[label=(\Alph*),leftmargin=*,itemsep=0.2em]
\item entrepreneurs
\item enthusiasts
\item miners
\item minions
\end{enumerate}
\par
\begingroup
\setlength{\parskip}{0pt}
\setlength{\parindent}{0pt}
\textbf{Review}\par
Answer: C\par
Confidence: high\par
Evidence refs: lines 19-20\par
Evidence quote: "Not all Forty-Niners, as the prospectors came to be known"\par
Jawaban (C) karena prospectors dalam passage merujuk pada orang yang mencari emas, yaitu miners.\par
\endgroup
\par
\vspace{0.25\baselineskip}
\noindent\textbf{20.} It can be inferred from the passage that many Forty Niners
\begin{enumerate}[label=(\Alph*),leftmargin=*,itemsep=0.2em]
\item acquired great wealth
\item infected many laborers
\item attained great dispositions
\item inflated many recounts
\end{enumerate}
\par
\begingroup
\setlength{\parskip}{0pt}
\setlength{\parindent}{0pt}
\textbf{Review}\par
Answer: A\par
Confidence: medium\par
Evidence refs: lines 20-24\par
Evidence quote: "Not all Forty-Niners ... made their journey in vain ... yields of \$300 to \$400 a day were not infrequent"\par
Jawaban (A) karena jika tidak semua gagal dan hasil harian ratusan dolar cukup sering terjadi, masuk akal disimpulkan banyak Forty-Niners memang memperoleh kekayaan besar.\par
\endgroup
\par
\vspace{0.25\baselineskip}
\noindent\textbf{21.} The author of the passage implies that the California Gold Fever
\begin{enumerate}[label=(\Alph*),leftmargin=*,itemsep=0.2em]
\item was a unique event in North American history
\item was a singular event in the succession of others
\item presented a lone occurrence on the North American continent
\item preserved a unique page in North American development
\end{enumerate}
\par
\begingroup
\setlength{\parskip}{0pt}
\setlength{\parindent}{0pt}
\textbf{Review}\par
Answer: B\par
Confidence: high\par
Evidence refs: lines 29-30\par
Evidence quote: "The gold rush of 1849 was followed by the next wave of migrants headed to the Klondike region"\par
Jawaban (B) karena penulis secara implisit menyatakan California Gold Fever bukan peristiwa tunggal yang sepenuhnya unik, melainkan satu peristiwa yang kemudian diikuti gelombang gold rush lain.\par
\endgroup
\par
\vspace{0.25\baselineskip}
\subsection*{Questions 22-31}
\begingroup
\small
\sloppy
\setlength{\parskip}{0pt}
\setlength{\parindent}{0pt}
During the late Middle Ages, oil paint took hold as the artistic medium of choice because\par
it was effective, flexible, and resilient relative to the wax-based, watercolor, fresco, or tem-\par
pera paints prevalent at the time. Although contemporary commercially prepared paints\par
contain a mixture of pigments and linseed oil, poppy oil paints are also available to con-\par
noisseurs. The original recipes developed in medieval European monasteries relied on fast-\par
drying bases derived from various organic oils predominantly valued for their medicinal\par
qualities. The pigments are insoluble, lightproof, and chemically inert powders ground in\par
the base. Occasionally, varnish can be added to increase the paste's ability to reflect light\par
and to cover pictures with a protective seal. The resulting stiff, resinous compounds are\par
often packaged in flexible metal or plastic tubes. Historically, yellow pigments have been\par
added to the oil, and then the paste was layered over tin foil to imitate the appearance of\par
gold leaf.\par
Despite the numerous experiments to accelerate the drying process, oil paints dry com-\par
paratively slowly. with little color alteration. An important advantage of color stability is that\par
tones and undertones are easy to blend, match, transpose, and grade, and mistakes and\par
smudges are simple to correct. Due to the creamy consistency of most mixtures, artists can\par
exploit their viscosity in thick applications, sprays, thin trickles, and three-dimensional\par
blobs. The purification by boiling and filtering and bleaching of oils can impart varied hues\par
to powdered pigments, while drying time can be reduced by adding metallic oxides.\par
Professional painters who mix their own medium usually have their own trademark meth-\par
ods of mixing materials that art experts recognize as a part of an artist's creative work.\par
The thickness of the paste also plays an important role in defining the stages of painting\par
a picture. After the basic design is sketched in pencil or charcoal, the broad background or\par
foreground areas of the canvas are covered with thin, diluted paint on top of the primer. A\par
thicker paint, often with added varnish, is subsequently used to refine and outline the foun-\par
dation. The width of the brush depends on the type of paint the artist chooses to use,\par
and stiff bristles are usually found in narrow brushes for making sharp lines, while softer\par
brushes of animal hair can be employed in broad strokes.\par
\endgroup
\medskip
\noindent\textbf{22.} What does the passage mainly discuss?
\begin{enumerate}[label=(\Alph*),leftmargin=*,itemsep=0.2em]
\item The evolution and history of oil paintings and media
\item The technology and development of drying oils
\item The recipes and ingredients for producing oil paints
\item The composition and techniques for mixing oil paints
\end{enumerate}
\par
\begingroup
\setlength{\parskip}{0pt}
\setlength{\parindent}{0pt}
\textbf{Review}\par
Answer: D\par
Confidence: high\par
Evidence refs: lines 3-9, lines 13-28\par
Evidence quote: "contain a mixture of pigments and linseed oil... varnish can be added... oil paints dry comparatively slowly... artists can exploit their viscosity... The thickness of the paste also plays an important role"\par
Jawaban (D) karena passage terutama membahas apa saja komposisi oil paint dan bagaimana sifat campurannya dipakai dalam teknik melukis. Unsur sejarah hanya pembuka, bukan fokus utama.\par
\endgroup
\par
\vspace{0.25\baselineskip}
\noindent\textbf{23.} It can be inferred from the passage that oil paintings
\begin{enumerate}[label=(\Alph*),leftmargin=*,itemsep=0.2em]
\item supplanted the use of tempera and fresco
\item took hold of the artistic choices in the Middle Ages
\item promoted artistic talent since the early times
\item supported the usefulness of applying paints
\end{enumerate}
\par
\begingroup
\setlength{\parskip}{0pt}
\setlength{\parindent}{0pt}
\textbf{Review}\par
Answer: A\par
Confidence: high\par
Evidence refs: lines 1-2\par
Evidence quote: "oil paint took hold as the artistic medium of choice... relative to the wax-based, watercolor, fresco, or tempera paints prevalent at the time"\par
Jawaban (A) karena jika oil paint menjadi `medium of choice` dibanding fresco dan tempera yang sebelumnya umum, implikasinya oil painting menggantikan dominasi media-media itu.\par
\endgroup
\par
\vspace{0.25\baselineskip}
\noindent\textbf{24.} In lines 4 and 5, the word "connoisseurs" is closest in meaning to
\begin{enumerate}[label=(\Alph*),leftmargin=*,itemsep=0.2em]
\item explorers
\item experts
\item exporters
\item experimenters
\end{enumerate}
\par
\begingroup
\setlength{\parskip}{0pt}
\setlength{\parindent}{0pt}
\textbf{Review}\par
Answer: B\par
Confidence: high\par
Evidence refs: lines 4-5\par
Evidence quote: "poppy oil paints are also available to connoisseurs"\par
Jawaban (B) karena `connoisseurs` berarti orang yang punya pengetahuan dan apresiasi mendalam, paling dekat dengan `experts`.\par
\endgroup
\par
\vspace{0.25\baselineskip}
\noindent\textbf{25.} According to the passage, medieval monks extracted oil
\begin{enumerate}[label=(\Alph*),leftmargin=*,itemsep=0.2em]
\item from minerals
\item in conjunction with pigments
\item from plants
\item in combination with medicines
\end{enumerate}
\par
\begingroup
\setlength{\parskip}{0pt}
\setlength{\parindent}{0pt}
\textbf{Review}\par
Answer: C\par
Confidence: medium\par
Evidence refs: lines 5-6\par
Evidence quote: "relied on fast-drying bases derived from various organic oils"\par
Jawaban (C) karena oil base disebut berasal dari `organic oils`, yang paling logis berarti sumber hayati seperti tanaman, bukan mineral atau obat. Namun teks tidak menyebut jenis tanamannya secara eksplisit.\par
Needs manual review: true\par
\endgroup
\par
\vspace{0.25\baselineskip}
\noindent\textbf{26.} In line 8, the phrase "the base" is closest in meaning to
\begin{enumerate}[label=(\Alph*),leftmargin=*,itemsep=0.2em]
\item paint
\item oil
\item chemicals.
\item pestle
\end{enumerate}
\par
\begingroup
\setlength{\parskip}{0pt}
\setlength{\parindent}{0pt}
\textbf{Review}\par
Answer: B\par
Confidence: high\par
Evidence refs: lines 5-8\par
Evidence quote: "bases derived from various organic oils... The pigments are... powders ground in the base"\par
Jawaban (B) karena `the base` merujuk ke oil base yang disebut tepat sebelumnya, yaitu dasar cair atau medium berbahan minyak.\par
\endgroup
\par
\vspace{0.25\baselineskip}
\noindent\textbf{27.} The purpose of paragraph 2 is to illustrate
\begin{enumerate}[label=(\Alph*),leftmargin=*,itemsep=0.2em]
\item the laboriousness of making oil paints
\item the durability of oil colors
\item the complexity of oil purification
\item the superiority of oil paints
\end{enumerate}
\par
\begingroup
\setlength{\parskip}{0pt}
\setlength{\parindent}{0pt}
\textbf{Review}\par
Answer: D\par
Confidence: high\par
Evidence refs: lines 13-19\par
Evidence quote: "oil paints dry comparatively slowly, with little color alteration... tones and undertones are easy to blend... mistakes and smudges are simple to correct"\par
Jawaban (D) karena paragraf 2 memaparkan kelebihan oil paint: warna stabil, mudah dibaurkan, mudah dikoreksi, dan bisa dipakai dengan berbagai tekstur. Itu menunjukkan superioritasnya, bukan sekadar daya tahan.\par
\endgroup
\par
\vspace{0.25\baselineskip}
\noindent\textbf{28.} In line 17, the word "viscosity" is closest in meaning to
\begin{enumerate}[label=(\Alph*),leftmargin=*,itemsep=0.2em]
\item stiffness
\item elasticity
\item stickiness
\item eloquence
\end{enumerate}
\par
\begingroup
\setlength{\parskip}{0pt}
\setlength{\parindent}{0pt}
\textbf{Review}\par
Answer: A\par
Confidence: medium\par
Evidence refs: lines 16-17\par
Evidence quote: "Due to the creamy consistency of most mixtures, artists can exploit their viscosity in thick applications"\par
Jawaban (A) karena dalam konteks ini `viscosity` menunjuk pada tingkat kekentalan atau ketegaran aliran cat, sehingga paling dekat dengan `stiffness`. Pilihan ini tidak sempurna, tetapi lebih dekat daripada `stickiness`.\par
Needs manual review: true\par
\endgroup
\par
\vspace{0.25\baselineskip}
\noindent\textbf{29.} Which of the following is NOT mentioned as components of oil paints?
\begin{enumerate}[label=(\Alph*),leftmargin=*,itemsep=0.2em]
\item accelerants
\item retarders
\item sealants
\item glosses
\end{enumerate}
\par
\begingroup
\setlength{\parskip}{0pt}
\setlength{\parindent}{0pt}
\textbf{Review}\par
Answer: B\par
Confidence: high\par
Evidence refs: lines 3-9, lines 18-19\par
Evidence quote: "pigments and linseed oil... varnish can be added... protective seal; drying time can be reduced by adding metallic oxides"\par
Jawaban (B) karena passage menyebut gloss atau sealant melalui varnish dan accelerant melalui metallic oxides, tetapi tidak pernah menyebut retarders yang memperlambat pengeringan.\par
\endgroup
\par
\vspace{0.25\baselineskip}
\noindent\textbf{30.} In line 20, the word "trademark" is closest in meaning to
\begin{enumerate}[label=(\Alph*),leftmargin=*,itemsep=0.2em]
\item signature
\item selection
\item significance
\item secret
\end{enumerate}
\par
\begingroup
\setlength{\parskip}{0pt}
\setlength{\parindent}{0pt}
\textbf{Review}\par
Answer: A\par
Confidence: high\par
Evidence refs: lines 20-21\par
Evidence quote: "their own trademark methods of mixing materials that art experts recognize as a part of an artist's creative work"\par
Jawaban (A) karena `trademark methods` di sini berarti cara khas yang menjadi ciri pengenal seniman, jadi paling dekat dengan `signature`.\par
\endgroup
\par
\vspace{0.25\baselineskip}
\noindent\textbf{31.} The author of the passage implies that an oil painting
\begin{enumerate}[label=(\Alph*),leftmargin=*,itemsep=0.2em]
\item requires professional painters to mix their own paint
\item contains a layer of canvas and charcoal
\item thickens as the oil continues to dry in stages
\item requires multiple layers of brushwork
\end{enumerate}
\par
\begingroup
\setlength{\parskip}{0pt}
\setlength{\parindent}{0pt}
\textbf{Review}\par
Answer: D\par
Confidence: high\par
Evidence refs: lines 22-28\par
Evidence quote: "After the basic design is sketched... areas of the canvas are covered with thin, diluted paint... A thicker paint... is subsequently used to refine and outline"\par
Jawaban (D) karena penulis menjelaskan proses bertahap: sketsa, lapisan tipis awal, lalu lapisan lebih tebal untuk penyempurnaan. Itu mengimplikasikan oil painting membutuhkan beberapa lapis pekerjaan kuas.\par
\endgroup
\par
\vspace{0.25\baselineskip}
\subsection*{Questions 32-41}
\begingroup
\small
\sloppy
\setlength{\parskip}{0pt}
\setlength{\parindent}{0pt}
Samuel Johnson's achievements in lexicography and writing have few rivals in English\par
eighteenth-century literature. Johnson was born on September 18, 1709, in the family of a\par
bookseller and stationer, and his education began informally among his father's books, as\par
Samuel was above all an insatiable reader. During his years at Pembroke College, Oxford,\par
he acquired a reputation for prodigious learning and eccentricities that were as diverse as\par
his reading. After being plagued by many ailments and emotional breakdowns, he was\par
compelled to leave the college in poverty, and later in life, his despondency led him to a\par
profound devotion to religion.\par
Following his father's death in 1731, Johnson tried his hand in teaching and attempted to\par
run a school, but his early ventures met with rapid demise. When Johnson was twenty-six,\par
he married Elizabeth Porter, a widow with a little fortune, who was twenty years his senior\par
and who became the source of calm, support, and self-confidence in his life. In 1737 in\par
London, Johnson began a life-long period of writing and lexicographic work that made him\par
a figure of renown.\par
To his contemporaries, Johnson was known as a man of large and unwieldy constitution\par
and unpleasant manner. He recklessly dominated conversations, monopolized discussion\par
topics, and boorishly overwhelmed opposing points of view. His antagonistic habits cost\par
him many a friend. In 1747, burdened with extreme poverty, Johnson was commissioned to\par
compile a dictionary by a consortium of booksellers who sought to raise their volume of\par
sales. After eight years of obsessive and consuming labor, this monumental work with\par
40,000 entries with lucid, witty, occasionally spiteful, and idiosyncratic definitions, titled\par
Dictionary of the English Language, saw the light of day in 1755. The remarkable, still-\par
quoted lexicographic work with an enormous range of elucidating descriptions and illustrative\par
exemplars became the imitable forerunner of how English language dictionaries are con-\par
structed to this day.\par
During the nineteenth century, literary critics largely saw Johnson's literary contributions\par
as slight with regard to their scholarly influence. It was perhaps Johnson's personal attrib-\par
utes that influenced the perspective through which he was portrayed. While his prose was\par
viewed as dignified, vivid, and urbane, his vast body of work in poetry, novels, essays,\par
tragedies, travelogues, biographies, diaries, and satire remained obscure until Joseph Krutch,\par
an American critic, published a penetrating study Samuel Johnson in 1944. In 1978, Walter\par
Bate won a Pulitzer Prize for his comprehensive psychological biography of Johnson that\par
established him as rational, compassionate, and undaunted advocate of moral order and\par
unceasing quest for truth.\par
\endgroup
\medskip
\noindent\textbf{32.} What does the passage mainly discuss?
\begin{enumerate}[label=(\Alph*),leftmargin=*,itemsep=0.2em]
\item The life, ailments, and dictionary of Samuel Johnson
\item Samuel Johnson's controversial character and work
\item Samuel Johnson's contradictory traits and publications
\item The family, commissions, and biographies of Samuel Johnson
\end{enumerate}
\par
\begingroup
\setlength{\parskip}{0pt}
\setlength{\parindent}{0pt}
\textbf{Review}\par
Answer: A\par
Confidence: high\par
Evidence refs: lines 1-2, lines 13-14, lines 18-25\par
Evidence quote: "Samuel Johnson's achievements in lexicography and writing have few rivals ... Johnson began a life-long period of writing and lexicographic work ... Dictionary of the English Language ... became the imitable forerunner of how English language dictionaries are constructed to this day."\par
Jawaban (A) karena passage terutama memberi ringkasan biografis tentang kehidupan Johnson, penyakit dan kesulitannya, lalu pencapaian besarnya sebagai leksikografer lewat kamusnya. Opsi lain hanya menyorot sebagian aspek.\par
\endgroup
\par
\vspace{0.25\baselineskip}
\noindent\textbf{33.} In line 5, the word prodigious is closest in meaning to
\begin{enumerate}[label=(\Alph*),leftmargin=*,itemsep=0.2em]
\item immoral
\item imminent
\item immense
\item immediate
\end{enumerate}
\par
\begingroup
\setlength{\parskip}{0pt}
\setlength{\parindent}{0pt}
\textbf{Review}\par
Answer: C\par
Confidence: high\par
Evidence refs: line 5\par
Evidence quote: "he acquired a reputation for prodigious learning"\par
Jawaban (C) karena prodigious dalam konteks learning berarti sangat besar atau luar biasa besarnya, paling dekat dengan immense.\par
\endgroup
\par
\vspace{0.25\baselineskip}
\noindent\textbf{34.} In line 6, the word "plagued" is closest in meaning to
\begin{enumerate}[label=(\Alph*),leftmargin=*,itemsep=0.2em]
\item diseased
\item disfigured
\item trampled
\item tormented
\end{enumerate}
\par
\begingroup
\setlength{\parskip}{0pt}
\setlength{\parindent}{0pt}
\textbf{Review}\par
Answer: D\par
Confidence: high\par
Evidence refs: lines 6-7\par
Evidence quote: "After being plagued by many ailments and emotional breakdowns"\par
Jawaban (D) karena plagued by ailments and emotional breakdowns berarti terus-menerus diganggu atau disiksa oleh masalah itu, jadi paling dekat dengan tormented.\par
\endgroup
\par
\vspace{0.25\baselineskip}
\noindent\textbf{35.} In line 4, the word "overlap" is closest in meaning to
\begin{enumerate}[label=(\Alph*),leftmargin=*,itemsep=0.2em]
\item coincide
\item consign
\item collide
\item concur
\end{enumerate}
\par
\begingroup
\setlength{\parskip}{0pt}
\setlength{\parindent}{0pt}
\textbf{Review}\par
Answer: A\par
Confidence: low\par
Evidence refs: line 4\par
Evidence quote: "Samuel was above all an insatiable reader."\par
Saya usulkan (A) hanya jika item ini memang bermaksud menguji kata overlap secara umum, karena overlap paling dekat dengan coincide. Namun source passage pada line 4 tidak memuat kata overlap sama sekali, jadi item ini cacat sumber dan perlu review manual.\par
Needs manual review: true\par
\endgroup
\par
\vspace{0.25\baselineskip}
\noindent\textbf{36.} According to the passage, Johnson's marriage
\begin{enumerate}[label=(\Alph*),leftmargin=*,itemsep=0.2em]
\item was a miserable failure due to poverty
\item was motivated by convenience and money
\item provided him a safe haven
\item afforded him new writing venues
\end{enumerate}
\par
\begingroup
\setlength{\parskip}{0pt}
\setlength{\parindent}{0pt}
\textbf{Review}\par
Answer: C\par
Confidence: high\par
Evidence refs: lines 11-12\par
Evidence quote: "he married Elizabeth Porter ... who became the source of calm, support, and self-confidence in his life"\par
Jawaban (C) karena passage menyatakan langsung bahwa istrinya menjadi sumber ketenangan, dukungan, dan rasa percaya diri, yaitu semacam safe haven bagi Johnson.\par
\endgroup
\par
\vspace{0.25\baselineskip}
\noindent\textbf{37.} The word "antagonistic" in line 17 is closest in meaning to.
\begin{enumerate}[label=(\Alph*),leftmargin=*,itemsep=0.2em]
\item hallow
\item heated
\item hospitable
\item hostile
\end{enumerate}
\par
\begingroup
\setlength{\parskip}{0pt}
\setlength{\parindent}{0pt}
\textbf{Review}\par
Answer: D\par
Confidence: high\par
Evidence refs: lines 16-18\par
Evidence quote: "He recklessly dominated conversations ... His antagonistic habits cost him many a friend."\par
Jawaban (D) karena antagonistic dijelaskan lewat perilaku agresif dan merugikan hubungan pertemanan, sehingga maknanya paling dekat dengan hostile.\par
\endgroup
\par
\vspace{0.25\baselineskip}
\noindent\textbf{38.} According to the passage, the Dictionary descriptions of word meanings
\begin{enumerate}[label=(\Alph*),leftmargin=*,itemsep=0.2em]
\item are emulated by modern lexicographers
\item were not without peculiarities
\item are considered exemplary by linguists
\item were not consistently illustrated
\end{enumerate}
\par
\begingroup
\setlength{\parskip}{0pt}
\setlength{\parindent}{0pt}
\textbf{Review}\par
Answer: B\par
Confidence: high\par
Evidence refs: lines 20-24\par
Evidence quote: "40,000 entries with lucid, witty, occasionally spiteful, and idiosyncratic definitions ..."\par
Jawaban (B) karena passage menekankan bahwa definisinya bukan hanya jelas dan cerdas, tetapi juga kadang spiteful dan idiosyncratic. Itu menunjukkan deskripsi makna katanya memang tidak tanpa keanehan atau peculiarities.\par
\endgroup
\par
\vspace{0.25\baselineskip}
\noindent\textbf{39.} The word "obscure" in line 30 is closest in meaning to
\begin{enumerate}[label=(\Alph*),leftmargin=*,itemsep=0.2em]
\item offensive
\item obtuse
\item irrelevant
\item irreverent
\end{enumerate}
\par
\begingroup
\setlength{\parskip}{0pt}
\setlength{\parindent}{0pt}
\textbf{Review}\par
Answer: C\par
Confidence: low\par
Evidence refs: lines 29-31\par
Evidence quote: "his vast body of work ... remained obscure until Joseph Krutch ... published a penetrating study"\par
Saya usulkan (C) sebagai pendekatan paling dekat, karena obscure di konteks ini berarti kurang dikenal atau kurang mendapat perhatian. Namun tidak ada opsi yang benar-benar setara dengan makna itu, jadi item ini juga perlu review manual.\par
Needs manual review: true\par
\endgroup
\par
\vspace{0.25\baselineskip}
\noindent\textbf{40.} Which of the following is NOT mentioned among the genres of Johnson's writing?
\begin{enumerate}[label=(\Alph*),leftmargin=*,itemsep=0.2em]
\item Plays
\item Biographies
\item Quotations
\item Articles
\end{enumerate}
\par
\begingroup
\setlength{\parskip}{0pt}
\setlength{\parindent}{0pt}
\textbf{Review}\par
Answer: C\par
Confidence: high\par
Evidence refs: lines 29-30\par
Evidence quote: "poetry, novels, essays, tragedies, travelogues, biographies, diaries, and satire"\par
Jawaban (C) karena quotations tidak disebut dalam daftar genre tulisan Johnson. Plays tercakup oleh tragedies, biographies disebut langsung, dan articles paling dekat dengan essays.\par
\endgroup
\par
\vspace{0.25\baselineskip}
\noindent\textbf{41.} The word "slight" in line 27 is closest in meaning to
\begin{enumerate}[label=(\Alph*),leftmargin=*,itemsep=0.2em]
\item trivial
\item triumphant
\item sinister
\item solemn
\end{enumerate}
\par
\begingroup
\setlength{\parskip}{0pt}
\setlength{\parindent}{0pt}
\textbf{Review}\par
Answer: A\par
Confidence: high\par
Evidence refs: lines 26-27\par
Evidence quote: "literary critics largely saw Johnson's literary contributions as slight with regard to their scholarly influence"\par
Jawaban (A) karena slight di sini berarti kecil atau tidak penting pengaruhnya, paling dekat dengan trivial.\par
\endgroup
\par
\vspace{0.25\baselineskip}
\subsection*{Questions 42-50}
\begingroup
\small
\sloppy
\setlength{\parskip}{0pt}
\setlength{\parindent}{0pt}
Psychologists who work on motivation research a wide range of human traits and physi-\par
ological characteristics that include the effects of hunger, reward, and punishment, as well\par
as desires for power, tangible achievement, social acceptance, belongingness, self-esteem,\par
and self-actualization. A plethora of hypotheses developed in the nineteenth and twentieth\par
centuries have the goal of identifying causes of an organism's behavior that can be both\par
conscious and unconscious. The hierarchical organization of human needs is a theoretical\par
model, originally established by an American psychologist, Abraham Maslow, in 1954. The\par
needs located at the bottom of the pyramid are the essentials of physiological survival that\par
encompass oxygen, water, nutrition, rest, and avoidance of pain. Maslow's theory, grounded\par
in research, also stipulated that these are variable and, at least to some extent, may explain,\par
for example, food gratification. The second tier is rooted in the human need for safety, sta-\par
bility, and protection.\par
In the human life cycle, the needs for belonging are manifested in the desires to marry,\par
have a family, belong in a community or among similarly minded people. In part, the need\par
to belong can also show up in a search for particular types of occupations or careers. The\par
next level of the hierarchy in effect deals with two substrata, where the first presumes the\par
need for status, prestige, recognition, appreciation, and dominance, and the higher division\par
includes a conglomeration of emotionally centered traits that pivot on competence, confi-\par
dence, mastery, achievement, independence, and freedom.\par
The top tier is different from all others, and Maslow referred to it as growth motivation\par
and self-actualization. At the highest level, individuals seek to realize and put to use their\par
creativity, talent, leadership, curiosity, and understanding. At this level people can reach\par
their full potentials and accurately perceive and accept reality, seek privacy and depth in\par
personal relationships, resist enculturation, and develop social interests, compassion, and\par
humanity. In many cases, self-actualizers do not lead ordinary lives, choose growth over\par
safety, and cultivate peak experiences that leave their mark and change one for the better.\par
\endgroup
\medskip
\noindent\textbf{42.} According to the passage, what does psychology of motivation attempt to uncover?
\begin{enumerate}[label=(\Alph*),leftmargin=*,itemsep=0.2em]
\item Reactions to hunger and desires for power
\item Developments of human physiology and mind
\item Reasons for human conduct and moving forces
\item Organization of human society and hierarchies
\end{enumerate}
\par
\begingroup
\setlength{\parskip}{0pt}
\setlength{\parindent}{0pt}
\textbf{Review}\par
Answer: C\par
Confidence: high\par
Evidence refs: lines 1-6\par
Evidence quote: "Psychologists who work on motivation research a wide range of human traits ... hypotheses ... have the goal of identifying causes of an organism's behavior"\par
Jawaban (C) karena passage menjelaskan bahwa psikologi motivasi meneliti trait manusia dan menyusun hipotesis untuk mengidentifikasi penyebab perilaku; itu paling dekat dengan alasan perilaku dan daya pendorongnya.\par
\endgroup
\par
\vspace{0.25\baselineskip}
\noindent\textbf{43.} The word "plethora" in line 4 is closest in meaning to
\begin{enumerate}[label=(\Alph*),leftmargin=*,itemsep=0.2em]
\item oversimplification
\item overlap
\item overreach
\item overabundance
\end{enumerate}
\par
\begingroup
\setlength{\parskip}{0pt}
\setlength{\parindent}{0pt}
\textbf{Review}\par
Answer: D\par
Confidence: high\par
Evidence refs: line 4\par
Evidence quote: "A plethora of hypotheses developed in the nineteenth and twentieth centuries"\par
Jawaban (D) karena `plethora` dalam konteks ini berarti jumlah yang sangat banyak, paling dekat dengan `overabundance`.\par
\endgroup
\par
\vspace{0.25\baselineskip}
\noindent\textbf{44.} It can be inferred from the passage that in Maslow's hierarchy
\begin{enumerate}[label=(\Alph*),leftmargin=*,itemsep=0.2em]
\item the first layer of needs dominates other tiers
\item the highest level of the model supercedes the lower levels
\item the second layer of needs is more urgent than the first
\item the third level of the model is embedded in the fourth
\end{enumerate}
\par
\begingroup
\setlength{\parskip}{0pt}
\setlength{\parindent}{0pt}
\textbf{Review}\par
Answer: A\par
Confidence: medium\par
Evidence refs: lines 8-12\par
Evidence quote: "The needs located at the bottom of the pyramid are the essentials of physiological survival ... The second tier is rooted in the human need for safety, stability, and protection."\par
Jawaban (A) paling didukung karena lapisan pertama disebut sebagai kebutuhan esensial untuk survival fisiologis, sehingga secara inferensial ia lebih mendasar dan mendominasi urgensi tier di atasnya.\par
\endgroup
\par
\vspace{0.25\baselineskip}
\noindent\textbf{45.} In line 6, the word "trendy" is closest in meaning to
\begin{enumerate}[label=(\Alph*),leftmargin=*,itemsep=0.2em]
\item tender
\item experimental
\item fashionable
\item trusted
\end{enumerate}
\par
\begingroup
\setlength{\parskip}{0pt}
\setlength{\parindent}{0pt}
\textbf{Review}\par
Answer: C\par
Confidence: low\par
Evidence refs: line 6\par
Evidence quote: "conscious and unconscious."\par
Jawaban (C) dipilih dari makna leksikal `trendy` = `fashionable`, tetapi kata `trendy` tidak muncul pada line\_map passage ini sehingga item tampak anomali dan perlu cek manual.\par
Needs manual review: true\par
\endgroup
\par
\vspace{0.25\baselineskip}
\noindent\textbf{46.} The word "gratification" in line 11 is closest in meaning to
\begin{enumerate}[label=(\Alph*),leftmargin=*,itemsep=0.2em]
\item gratitude
\item pleasure
\item gratuity
\item pleasantry
\end{enumerate}
\par
\begingroup
\setlength{\parskip}{0pt}
\setlength{\parindent}{0pt}
\textbf{Review}\par
Answer: B\par
Confidence: high\par
Evidence refs: lines 10-11\par
Evidence quote: "may explain, for example, food gratification"\par
Jawaban (B) karena `gratification` pada frasa `food gratification` berarti rasa puas atau nikmat, paling dekat dengan `pleasure`.\par
\endgroup
\par
\vspace{0.25\baselineskip}
\noindent\textbf{47.} It can be inferred from the passage that in modern-day terms, the second layer of needs can be reflected in people's desire for
\begin{enumerate}[label=(\Alph*),leftmargin=*,itemsep=0.2em]
\item a house in an upscale neighborhood
\item a protected existence and dependence
\item a measure of job and financial security
\item a degree of friendship and family life
\end{enumerate}
\par
\begingroup
\setlength{\parskip}{0pt}
\setlength{\parindent}{0pt}
\textbf{Review}\par
Answer: C\par
Confidence: high\par
Evidence refs: lines 11-12\par
Evidence quote: "The second tier is rooted in the human need for safety, stability, and protection."\par
Jawaban (C) karena dalam istilah modern kebutuhan akan safety, stability, dan protection paling wajar tercermin pada keamanan kerja dan finansial.\par
\endgroup
\par
\vspace{0.25\baselineskip}
\noindent\textbf{48.} The word "conglomeration" in line 18 is closest in meaning to
\begin{enumerate}[label=(\Alph*),leftmargin=*,itemsep=0.2em]
\item complex
\item congress
\item conjunction
\item connotation
\end{enumerate}
\par
\begingroup
\setlength{\parskip}{0pt}
\setlength{\parindent}{0pt}
\textbf{Review}\par
Answer: A\par
Confidence: medium\par
Evidence refs: lines 18-19\par
Evidence quote: "includes a conglomeration of emotionally centered traits that pivot on competence, confidence, mastery, achievement, independence, and freedom"\par
Jawaban (A) karena `conglomeration` berarti kumpulan atau gabungan banyak unsur; dari opsi yang ada, `complex` paling dekat maknanya.\par
\endgroup
\par
\vspace{0.25\baselineskip}
\noindent\textbf{49.} Which of the following is NOT mentioned as a factor in human motivation?
\begin{enumerate}[label=(\Alph*),leftmargin=*,itemsep=0.2em]
\item Esteem
\item Participation
\item Accomplishment
\item Conformism
\end{enumerate}
\par
\begingroup
\setlength{\parskip}{0pt}
\setlength{\parindent}{0pt}
\textbf{Review}\par
Answer: D\par
Confidence: medium\par
Evidence refs: lines 1-4, lines 23-24\par
Evidence quote: "desires for power, tangible achievement, social acceptance, belongingness, self-esteem, and self-actualization ... resist enculturation"\par
Jawaban (D) karena esteem dan achievement disebut eksplisit, sedangkan `conformism` tidak disebut sebagai faktor motivasi; bahkan pada level tertinggi justru disebut `resist enculturation`.\par
\endgroup
\par
\vspace{0.25\baselineskip}
\noindent\textbf{50.} Which of the following conclusions is supported by the passage?
\begin{enumerate}[label=(\Alph*),leftmargin=*,itemsep=0.2em]
\item Genuine self-actualizers may attain self-satisfaction.
\item Sincere self-promoters can achieve full contentment.
\item Real self-starters can achieve their lives' desires.
\item True self-actualizers may lead complicated lives.
\end{enumerate}
\par
\begingroup
\setlength{\parskip}{0pt}
\setlength{\parindent}{0pt}
\textbf{Review}\par
Answer: D\par
Confidence: high\par
Evidence refs: lines 25-26\par
Evidence quote: "In many cases, self-actualizers do not lead ordinary lives, choose growth over safety, and cultivate peak experiences"\par
Jawaban (D) karena passage menyimpulkan bahwa self-actualizers sering tidak menjalani hidup yang biasa, sehingga kesimpulan paling didukung adalah mereka dapat menjalani hidup yang rumit atau tidak sederhana.\par
\endgroup
\par
\vspace{0.25\baselineskip}
\newpage
\section*{Answer Key Appendix}
\subsection*{Listening}
\begin{itemize}[leftmargin=*]
\item 1: C
\item 2: D
\item 3: D
\item 4: A
\item 5: B
\item 6: D
\item 7: A
\item 8: C
\item 9: B
\item 10: A
\item 11: D
\item 12: C
\item 13: A
\item 14: B
\item 15: C
\item 16: A
\item 17: D
\item 18: A
\item 19: C
\item 20: B
\item 21: B
\item 22: C
\item 23: A
\item 24: B
\item 25: C
\item 26: A
\item 27: C
\item 28: B
\item 29: A
\item 30: B
\item 31: B
\item 32: D
\item 33: B
\item 34: C
\item 35: B
\item 36: A
\item 37: C
\item 38: D
\item 39: C
\item 40: B
\item 41: C
\item 42: C
\item 43: B
\item 44: D
\item 45: D
\item 46: B
\item 47: C
\item 48: C
\item 49: D
\item 50: B
\end{itemize}
\subsection*{Structure}
\begin{itemize}[leftmargin=*]
\item 1: A
\item 2: C
\item 3: D
\item 4: A
\item 5: B
\item 6: C
\item 7: A
\item 8: C
\item 9: C
\item 10: D
\item 11: C
\item 12: A
\item 13: C
\item 14: D
\item 15: C
\item 16: A
\item 17: B
\item 18: C
\item 19: B
\item 20: D
\item 21: B
\item 22: A
\item 23: D
\item 24: B
\item 25: A
\item 26: D
\item 27: B
\item 28: D
\item 29: C
\item 30: B
\item 31: A
\item 32: C
\item 33: B
\item 34: B
\item 35: D
\item 36: C
\item 37: B
\item 38: B
\item 39: D
\item 40: D
\end{itemize}
\subsection*{Reading}
\begin{itemize}[leftmargin=*]
\item 1: C
\item 2: D
\item 3: A
\item 4: B
\item 5: B
\item 6: A
\item 7: B
\item 8: A
\item 9: B
\item 10: C
\item 11: B
\item 12: D
\item 13: B
\item 14: C
\item 15: A
\item 16: B
\item 17: C
\item 18: D
\item 19: C
\item 20: A
\item 21: B
\item 22: D
\item 23: A
\item 24: B
\item 25: C
\item 26: B
\item 27: D
\item 28: A
\item 29: B
\item 30: A
\item 31: D
\item 32: A
\item 33: C
\item 34: D
\item 35: A
\item 36: C
\item 37: D
\item 38: B
\item 39: C
\item 40: C
\item 41: A
\item 42: C
\item 43: D
\item 44: A
\item 45: C
\item 46: B
\item 47: C
\item 48: A
\item 49: D
\item 50: D
\end{itemize}
\end{document}
